% Appendix: design choices of the semi-synthetic data generation.
% Include after the existing appendices of second_draft.tex with
%     % Appendix: design choices of the semi-synthetic data generation.
% Include after the existing appendices of second_draft.tex with
%     % Appendix: design choices of the semi-synthetic data generation.
% Include after the existing appendices of second_draft.tex with
%     % Appendix: design choices of the semi-synthetic data generation.
% Include after the existing appendices of second_draft.tex with
%     \input{appendix_semisynthetic}
% Needs only amsmath and booktabs, both already loaded there. Self-contained: no \todo,
% no external figures. Code references are to src/CausalSurv/semisynthetic/.

\providecommand{\drv}[1]{\texttt{#1}}
\providecommand{\arm}[1]{\textsc{#1}}
\providecommand{\Ind}{\mathbf{1}}

\section{Semi-Synthetic Data Generation}
\label{app:semisynthetic}

This appendix records every design choice of the semi-synthetic benchmark used for
counterfactual validation (Section~\ref{sec:counterfactual}): what was decided, and why.
The benchmark keeps the real covariates and treatment history of the cohort, and
\emph{simulates} the treatment of the line being scored, the potential survival time under
every arm, and the censoring. The true survival curve, restricted mean survival time (RMST)
and best arm of every patient under every arm are then known in closed form.

The numbers quoted below are for the reference configuration (assignment strength
$\gamma=1$, hidden-confounder strength $s=0$, heterogeneity $h=1$). All coefficients are
set by hand, not estimated: they fix the \emph{difficulty} of the benchmark, not its realism.
DynaSurv's own results on the sweep (Section~\ref{app:semisynthetic:sweep}) are reported
against the oracle and naive Kaplan--Meier references only: this rebuild does not yet include
the Cox/DeepSurv baselines used for the observed-vs-hidden-confounding argument in
Section~\ref{sec:counterfactual}.

% -----------------------------------------------------------------------------------------
\subsection{Cohort and arms}
\label{app:semisynthetic:cohort}

\paragraph{Cohort.} HR$^+$/HER2$^-$ patients whose first treatment line starts in 2018 or
later, lines 1--4, built exactly as the training pipeline builds its cohort (patient-level
entry-year filter, inner join with the static file), so that simulated rows line up one to
one with the rows the model trains on. This gives $6{,}191$ patients and $13{,}116$
(patient, line) records, $6{,}191 / 3{,}401 / 2{,}169 / 1{,}355$ on lines 1 to 4. The year
restriction is the identifiability window of the real pipeline: every arm was available
throughout. One exact duplicate (patient, line) record in the source file is removed, because
after expansion it would become a sample of its own with a false ``previous line''.
Administrative censoring is measured from the last observed line start, 2024-02-21.

\paragraph{Arms.} Four arms: \arm{ET alone}, \arm{ET+anti-CDK wo CT}, \arm{MonoCT std alone},
\arm{PolyCT alone}. \arm{CT+anti-angio} is not simulated (21, 7, 2 and 2 records per line in the
restricted cohort: no overlap). It and all other categories (\arm{CT+anti-HER2}, \arm{CT+IT},
\arm{CT+TT}, \arm{ET+TT}, \arm{No treatment}, \arm{Other}) remain in the data as
\emph{history} only and are excluded from the set of recommendable arms. Arm indices follow
the alphabetical order of the arm names throughout.

\paragraph{Real history, simulated present.} Covariates, past treatments and line structure
(previous-line duration, gap, line number) are real. Only the arm and outcome of the line
being scored are simulated, and the simulated outcome does \emph{not} feed back into later
lines. This keeps the covariate distribution real and makes each line a separate problem
given its history. The price is that the benchmark does not test how a model handles
treatment--outcome feedback across lines.

% -----------------------------------------------------------------------------------------
\subsection{Prefix expansion}
\label{app:semisynthetic:expansion}

A patient with $L$ lines becomes $L$ samples. Sample $j\in\{1,\dots,L\}$ contains the real
records of lines $1,\dots,j-1$ and record $j$ with the simulated arm and simulated outcome;
its identifier is $10\cdot\text{id}+j$, and all samples of a patient share the original id
(the unit of any train/validation grouping) and the entry year. The reference cohort yields
$13{,}116$ samples in $24{,}920$ rows.

The reason is that the sequence model steps through lines in order, and its hidden state at
line $j$ does not depend on the arm at line $j$. Scoring only the last line of each sample
therefore requires no change to the model or its losses; the alternative, simulating all
lines of a patient jointly, would require outcome feedback and a joint model. The data
module masks every row except the sample's last line; history rows carry placeholder
outcomes ($0$).

Columns describing the real current line and its outcome (drug flags other than the treatment
category, all outcome columns other than the two survival targets, line end date, death date)
are removed, so that no real value sits next to a simulated one. Static features do not reach
the encoder, so age and menopausal status are copied into dynamic columns; otherwise they
would act as unplanned hidden confounders.

% -----------------------------------------------------------------------------------------
\subsection{Drivers and causal roles}
\label{app:semisynthetic:drivers}

Each sample has a vector $z\in\mathbb{R}^{16}$ of drivers computed from its real record.
Continuous drivers are standardised over the cohort; the others are $0/1$ flags.
Table~\ref{tab:semisynthetic:drivers} gives their definitions and roles. The role of each
driver is enforced when the configuration is loaded: a coefficient on an instrument in the
outcome equation, or on a prognostic factor in the assignment equation, is rejected.

\begin{table}[ht]
\centering\small
\caption{Drivers of the data-generating process. ``Both'' means the driver enters assignment
and outcome (confounder).}
\label{tab:semisynthetic:drivers}
\begin{tabular}{@{}lll@{}}
\toprule
Driver & Definition & Role \\
\midrule
\drv{liver} & liver metastases $>0$ & both \\
\drv{visceral} & pulmonary or pleural metastases $>0$ & both \\
\drv{bone\_only} & bone metastases, none of liver, visceral, brain & both \\
\drv{mpps} & imputed ECOG/WHO performance status, 0--4 (standardised) & both \\
\drv{log\_prev\_line} & $\log(1+\text{previous line duration})$ (standardised) & both \\
\drv{prev\_cdk}, \drv{prev\_et}, \drv{prev\_ct} & class of the \emph{real} previous arm (one-hot; all zero at line 1) & both \\
\drv{cum\_new\_sites} & cumulative new metastatic sites (standardised) & both \\
\drv{age}, \drv{menopause} & age at selection (standardised), menopausal status & both \\
\drv{calendar} & months since 2018-01-01 (standardised) & assignment only \\
\drv{brain} & brain metastases $>0$ & outcome only \\
\drv{lobular} & invasive lobular histology & outcome only \\
\drv{brca} & \textit{BRCA1} or \textit{BRCA2} & outcome only \\
\drv{old\_site\_progression} & progression of an old site $>0$ & outcome only \\
\bottomrule
\end{tabular}
\end{table}

The model still receives all real covariate columns ($73$, plus the calendar covariate); only
these $16$ matter for the simulated data, and the remainder are nuisance columns it must
learn to ignore. A buffer-time covariate is deliberately not a driver.

\paragraph{Hidden confounder.} A latent variable
$u=\rho\,\tilde z_{\text{liver}}+\sqrt{1-\rho^2}\,\varepsilon$ with $\varepsilon\sim\mathcal{N}(0,1)$,
$\tilde z_{\text{liver}}$ the standardised liver indicator and $\rho=0.3$, enters both
assignment and outcome and is \emph{never} written to the model input. Correlating it with a
real covariate keeps it from being independent noise. Its outcome coefficient is fixed and
only its coupling to assignment is swept, so that strength $0$ is an unconfounded reference
carrying the same outcome heterogeneity: any error growth above it is confounding bias.

% -----------------------------------------------------------------------------------------
\subsection{Treatment assignment}
\label{app:semisynthetic:assignment}

For a sample on line $\ell$ with drivers $z$ and hidden $u$, the arm $A$ is drawn with
\begin{equation}
    \pi_a(z,u)=\mathbb{P}(A=a\mid z,u)
    =\frac{\exp\eta^{\text{A}}_a}{\sum_{b}\exp\eta^{\text{A}}_b},
    \qquad
    \eta^{\text{A}}_a=\alpha_{\ell a}+\gamma\,B_a^{\top}z+s\,c_a\,u,
    \label{eq:semisynthetic:assignment}
\end{equation}
where $B_a$ are hand-set coefficients (Table~\ref{tab:semisynthetic:coefficients}, shared
across lines), $c_a$ the hidden loadings, $\gamma$ the assignment strength, and $s$ the hidden
strength. Coefficients are chosen for clinically plausible signs (visceral involvement pushes
towards combination chemotherapy, a long previous line towards endocrine-based therapy, a CDK
inhibitor is rarely repeated) and moderate size.

\paragraph{Calibrated intercepts.} The intercepts $\alpha_{\ell a}$ are not set by hand. For
every line they are fitted, by iterative proportional fitting
\begin{equation}
    \alpha_{\ell a}\leftarrow\alpha_{\ell a}+\log\frac{q_{\ell a}}{\bar\pi_{\ell a}},
    \qquad
    \bar\pi_{\ell a}=\frac{1}{n_\ell}\sum_{i\in\ell}\pi_a(z_i,u_i),
\end{equation}
until $\max_a|\bar\pi_{\ell a}-q_{\ell a}|<10^{-6}$, where $q_{\ell\cdot}$ is the real
distribution of the four arms on line $\ell$. The marginal arm mix therefore equals the real
mix at every $\gamma$ and $s$: a sweep changes \emph{who} receives which arm, never \emph{how
many}. The average runs over every sample of the line (each is given a simulated arm) while
the target is the real mix among the four arms only.
The targets $q_\ell$ (order: ET, ET+CDK, MonoCT, PolyCT) are
$(.183,.628,.090,.100)$, $(.148,.412,.330,.109)$, $(.049,.107,.611,.233)$ and
$(.032,.031,.580,.357)$ for lines 1 to 4.

\paragraph{Sampling.} The arm is drawn by inverting the cumulative distribution with
\emph{one uniform per sample}, so that draws coincide wherever propensities coincide. The true
propensities are stored.

\paragraph{Range of the sweep.} The effective sample size of inverse-propensity weights of the
assigned arm, divided by $n$, is $0.51$, $0.40$, $0.22$, $0.01$ and $0.00$ at $\gamma=0$, $0.5$,
$1$, $2$ and $4$ respectively (no hidden confounder), and $0.07$ at $s=2$ even with $\gamma=0$. Beyond $\gamma\approx1.5$ or
$s\approx1$ the benchmark measures extrapolation rather than confounding, so the sweep is
restricted to $\gamma\le1.5$ and $s\le1$.

\begin{table}[ht]
\centering\small
\caption{Hand-set coefficients (blank $=0$). Left: assignment coefficients $B_a$ and hidden
loadings $c_a$. Right: outcome prognostic term $f$ and arm interactions $g_a$. Continuous
drivers are per standard deviation, flags are on \emph{vs} off; outcome coefficients are on the
log-hazard scale (positive $=$ worse). The row ``Hidden'' gives the loadings $c_a$ of $u$ on the
arms' assignment logits and, in the $f$ column, its outcome coefficient $\beta_u$.}
\label{tab:semisynthetic:coefficients}
\setlength{\tabcolsep}{4pt}
\begin{tabular}{@{}l rrrr c rrrrr@{}}
\toprule
 & \multicolumn{4}{c}{Assignment $B_a$} & & \multicolumn{5}{c}{Outcome} \\
\cmidrule(lr){2-5}\cmidrule(l){7-11}
Driver & ET & CDK & Mono & Poly & & $f$ & $g_{\text{ET}}$ & $g_{\text{CDK}}$ & $g_{\text{Mono}}$ & $g_{\text{Poly}}$ \\
\midrule
\drv{liver}          &      & $-0.3$ & $0.5$  & $0.9$  & & $0.45$ & $0.30$ & $0.25$ & $-0.15$ & $-0.35$ \\
\drv{visceral}       &      & $-0.2$ & $0.3$  & $0.5$  & & $0.25$ &        &        &         & $-0.25$ \\
\drv{bone\_only}     &      & $0.5$  & $-0.4$ & $-0.6$ & & $-0.30$ & $-0.15$ & $-0.30$ &      &         \\
\drv{mpps}           &      & $-0.3$ & $0.3$  & $-0.2$ & & $0.30$ &        &        & $-0.10$ & $0.15$  \\
\drv{log\_prev\_line} &     & $0.4$  & $-0.4$ & $-0.5$ & & $-0.25$ &       & $-0.20$ &        &         \\
\drv{prev\_cdk}      & $-0.3$ & $-1.0$ & $0.4$ & $0.3$  & & $0.20$ &        & $0.40$ &         &         \\
\drv{prev\_et}       & $-0.5$ &      &        &        & &        & $0.30$ &        &         &         \\
\drv{prev\_ct}       &      &        & $0.3$  &        & &        &        &        &         &         \\
\drv{cum\_new\_sites} &     & $-0.1$ & $0.2$  & $0.3$  & & $0.15$ &        &        &         &         \\
\drv{age}            & $0.4$ &       & $0.3$  & $-0.4$ & & $0.10$ &        &        & $-0.10$ & $0.15$  \\
\drv{menopause}      &      & $0.3$  &        & $-0.3$ & & $0.05$ &        &        &         &         \\
\drv{calendar}       &      & $0.3$  & $-0.1$ & $-0.2$ & &        &        &        &         &         \\
\drv{brain}          &      &        &        &        & & $0.40$ &        &        &         &         \\
\drv{lobular}        &      &        &        &        & & $0.10$ &        &        &         &         \\
\drv{brca}           &      &        &        &        & & $0.15$ &        &        &         &         \\
\drv{old\_site\_prog.} &    &        &        &        & & $0.15$ &        &        &         &         \\
\midrule
Hidden ($c_a$; $\beta_u$ under $f$) & 0 & $-0.5$ & $0.5$ & $1.0$ & & $0.4$ & & & & \\
Arm effect $\tau_a$  &      &        &        &        & & & $0$ & $-0.20$ & $0.10$ & $0$ \\
\bottomrule
\end{tabular}
\end{table}

% -----------------------------------------------------------------------------------------
\subsection{Potential outcomes}
\label{app:semisynthetic:outcome}

Time from the start of the line to death under arm $a$ follows a Weibull model with
line-specific shape $k_\ell$ and scale $\lambda_\ell$,
\begin{equation}
    S_a(t\mid z,u)=\exp\!\Big(-\big(t/\lambda_\ell\big)^{k_\ell}e^{\eta_a}\Big),
    \qquad
    \eta_a=f(z)+\tau_a+h\,g_a(z)+\beta_u\,u,
    \label{eq:semisynthetic:outcome}
\end{equation}
that is, $T_a\sim\text{Weibull}(k_\ell,\lambda_\ell e^{-\eta_a/k_\ell})$; $\eta_a>0$ is a higher
hazard. $f$ is the prognostic term, $\tau_a$ the arm effect, $g_a$ the interaction that makes
the best arm differ across patients (who benefits from what), $h$ the heterogeneity, and
$\beta_u=0.4$ the fixed outcome coefficient of the hidden confounder. The functions $f$ and
$g_a$ are linear in $z$ (Table~\ref{tab:semisynthetic:coefficients}) and \emph{centred within
each line}, which makes $(k_\ell,\lambda_\ell)$ the baseline of the average patient.

\paragraph{Fitting the baseline.} For each line, $(k_\ell,\lambda_\ell)$ minimise the squared
difference between the population survival implied by the simulation for the assigned arms,
$\bar S_\ell(t)=n_\ell^{-1}\sum_{i\in\ell}\exp\!\big(-(t/\lambda)^{k}e^{\eta_{i,A_i}}\big)$, and the
real per-line Kaplan--Meier curve, over the observed jump times at which at least $5\%$ of the
line is still at risk. The result is $k_\ell\approx1.27$--$1.33$ and
$\lambda_\ell\approx53/27/22/18$ months; the fitted medians are $48.5/22.6/16.3/12.7$ months
against $47.7/22.8/16.5/12.4$ in the real data.

\paragraph{Closed-form truth.} Survival is given by \eqref{eq:semisynthetic:outcome}. With
$\sigma_a=\lambda_\ell e^{-\eta_a/k_\ell}$ the RMST to horizon $\tau$ is
\begin{equation}
    \text{RMST}_a(\tau)=\int_0^{\tau}S_a(t)\,dt
    =\sigma_a\,\Gamma\!\big(1+\tfrac1{k_\ell}\big)\,P\!\big(\tfrac1{k_\ell},(\tau/\sigma_a)^{k_\ell}\big),
\end{equation}
with $P$ the regularised lower incomplete gamma function; the individual treatment effect of
arm $a$ against $b$ is $\text{RMST}_a-\text{RMST}_b$ and the best arm maximises RMST over the
recommendable arms.

\paragraph{Coupling of the arms.} A single uniform per sample is shared by all arms, so
$T_a=\sigma_a(-\log V)^{1/k_\ell}$ with the same $V\sim\mathcal{U}(0,1)$ for every $a$. The
arms then differ only through $\eta_a$, and the individual effect is not swamped by independent
noise. The marginals are unaffected; the joint distribution of potential outcomes is one
particular choice, and another coupling would give the same marginals with different
individual-level effects.

\paragraph{Effect sizes.} Effects are small on the RMST scale: the average true pairwise
differences at the reference configuration lie between $-0.44$ and $+0.42$ months, and the
\emph{headroom} of a policy (value of the oracle minus value of the best single arm per line)
is $0.00, 0.11, 0.47, 1.25, 1.71$ months for $h=0,0.5,1,2,3$. The arm effect of
\arm{ET+anti-CDK} was first $-0.35$, which made it the best arm for $63$--$74\%$ of patients
(by line) so that ``always CDK'' scored highly; at $-0.20$ its share on line 2 drops from
$63\%$ to $48\%$ and the best-arm metric discriminates.

% -----------------------------------------------------------------------------------------
\subsection{Censoring}
\label{app:semisynthetic:censoring}

The observed time and event indicator are $\min(T,C)$ and $\Ind\{T\le C\}$ with
\begin{equation}
    C=\min\big(C^{\text{adm}},\,C^{\text{drop}}\big),\qquad
    C^{\text{adm}}=\frac{t_{\text{cutoff}}-t^{\text{start}}}{30.44},\qquad
    C^{\text{drop}}\sim\text{Exp}(r_\ell),
\end{equation}
in months, with $T$ the latent time under the simulated arm. Administrative censoring is a
fact of the real cohort and is kept sample by sample. Dropout is added only where the cutoff
alone leaves the simulated event rate above the real one; $r_\ell$ is found by bisection on
$\log r$ (60 steps) so that the event rate equals the real rate of the restricted cohort,
using fixed unit-exponential draws so that the event rate is a deterministic decreasing
function of $r$. The targets are the event rates of the restricted cohort,
$0.442/0.547/0.578/0.596$ (not those of the full 2008-onwards cohort, $0.73$--$0.85$). With
the cutoff alone the rates would be $0.465/0.587/0.639/0.679$, so dropout is active on every
line, at $0.0024/0.0041/0.0073/0.0103$ per month.

Dropout draws are independent conditional on line and the fitted generator parameters.
Administrative follow-up depends on calendar date, which also affects assignment and is
associated with patient characteristics. Marginal independent censoring is therefore not
guaranteed, and training-cohort censoring weights need not transfer to the later test cohort.
Protocol v2 uses known survival probabilities and uncensored latent outcomes for primary
simulation scoring instead of making that transfer assumption.

% -----------------------------------------------------------------------------------------
\subsection{Randomness and artefacts}
\label{app:semisynthetic:seeds}

Four independent random streams (hidden confounder, assignment, outcome, censoring) are
spawned from one seed sequence keyed by (data seed, replicate). A sweep over $\gamma$, $s$ or
$h$ therefore reuses the same $u$, the same assignment uniforms, the same latent uniforms and
the same dropout draws: levels differ \emph{only} by the swept knob, and repeating a
configuration reproduces the data exactly.

Each replicate writes (i) the expanded dynamic and re-keyed static tables under the names of
the real data files, so that the directory can be used in place of the real data; (ii) a truth
table with, per sample, the assigned arm, hidden $u$, Weibull parameters, latent time under
every arm, both censoring times, observed time and event, and propensity and $\eta$ for every
arm; and (iii) a manifest with every parameter, the calibrated intercepts, the arm counts per
line, the fitted Weibull parameters and medians, the dropout rates, and the achieved and
target event rates. The hidden $u$ appears only in the truth table.

% -----------------------------------------------------------------------------------------
\subsection{Training and data-loading choices}
\label{app:semisynthetic:training}

\begin{itemize}
    \item \textbf{Mask.} Only the last line of each sample is unmasked, and the counts that
    decide which arms are recommendable are computed on those rows.
    \item \textbf{Grouped splits.} Every split other than the temporal one (random hold-out,
    cross-validation folds, early-stopping split) groups by original patient, because the
    samples of a patient share their history and would otherwise leak across the boundary. A
    temporal split already keeps them together (shared entry year).
    \item \textbf{Protocol-v2 validation.} Twenty percent of the pre-2021 original patients
    form a separate development validation set (fixed split seed 0); the remainder train
    the model, scaler and support models. The survival grid is fitted from training endpoint
    outcomes only. The temporal test set never supplies early stopping or hyperparameter
    selection. Historical results below predate this correction and require rerunning.
    \item \textbf{Temporal hold-out.} Patients entering in 2021 or later ($26\%$ of samples),
    a deliberate policy-era shift; calendar time is an instrument of the assignment, so the
    model receives it.
    \item \textbf{Support thresholds.} Arms need at least $30$ samples, $10$ events and $30$
    samples followed to the horizon at a line to be recommendable, against $200/30/100$ on
    real data, because the four arms thin out on later lines (on line 4, \arm{ET alone} has
    $25$ samples and is excluded there).
    \item \textbf{Optimisation.} Early stopping on the validation loss (patience $20$, at most
    $150$ epochs). The architecture and learning-rate settings are those tuned on the real
    data; they were not re-tuned for the synthetic data.
\end{itemize}

% -----------------------------------------------------------------------------------------
\subsection{Evaluation against the truth}
\label{app:semisynthetic:evaluation}

All quantities are computed on the temporal hold-out, at each sample's own line $\ell$ and its
horizon $\tau_\ell$ ($24, 18, 12, 12$ months), on a $0.1$-month grid, and only for arms that
are recommendable at that line. Write $\hat S_a$ for a predicted curve and $S_a$ for the truth.
\begin{itemize}
    \item \textbf{Curve error.} $\big(\tau_\ell^{-1}\!\int_0^{\tau_\ell}(\hat S_a-S_a)^2dt\big)^{1/2}$
    averaged in squares over samples, and the error and bias of $\widehat{\text{RMST}}_a-\text{RMST}_a$,
    reported separately for the assigned (factual) arm and for the other (counterfactual) arms.
    \item \textbf{Effect error.} For each arm pair $(a,b)$,
    $\text{PEHE}=\big(\mathbb{E}[(\hat\Delta_{ab}-\Delta_{ab})^2]\big)^{1/2}$ with
    $\Delta_{ab}=\text{RMST}_a-\text{RMST}_b$, the bias of the average effect
    $\mathbb{E}[\hat\Delta_{ab}]-\mathbb{E}[\Delta_{ab}]$, and the share of samples with the
    correct sign. Hidden confounding can affect both bias and PEHE.
    \item \textbf{Policy.} A policy picks the arm with the highest predicted RMST among
    supported arms (abstaining to the assigned arm where none is supported). Value is the
    true RMST of the chosen arm; regret is the oracle value minus it; the hit rate is the share
    of samples where the chosen arm is the true best arm. Reference policies: the assigned
    arm, the best single arm per line (an upper bound for constant policies), and uniform
    random.
    \item \textbf{Factual accuracy (protocol v2).} At the exact requested horizon, write
    $p=S_A(\tau_\ell)$ and $q=\hat S_A(\tau_\ell)$. The primary Brier score averages
    $p(1-q)^2+(1-p)q^2$. We also report Brier scores from uncensored latent outcomes and
    concordance of latent event times with horizon-specific predicted risk. These are
    simulation diagnostics, not estimators available on real censored data. The historical
    IPCW results below used training-cohort censoring weights and must be replaced.
    \item \textbf{References.} The \emph{oracle} (true curves; must score zero error and zero
    regret) and the \emph{naive Kaplan--Meier} (per line and arm, no covariates; the bar a
    covariate model must clear).
\end{itemize}

% -----------------------------------------------------------------------------------------
\subsection{Sweep design}
\label{app:semisynthetic:sweep}

One knob is varied at a time, the others staying at the reference values.
Table~\ref{tab:semisynthetic:sweep} lists the cells; each is replicated three times (data seed
and training seed equal to the replicate index), giving $30$ runs. The hidden-confounder axis
is the one that separates methods that adjust for observed covariates only from methods that
would be unbiased under hidden confounding; the assignment axis is expected to be flat for
outcome-regression estimators, since observed confounding adds variance rather than bias.

\begin{table}[ht]
\centering\small
\caption{Sweep cells. The reference cell is $\gamma=1$, $s=0$, $h=1$; $s=0$ and $h=1$ are
therefore not repeated.}
\label{tab:semisynthetic:sweep}
\begin{tabular}{@{}lll@{}}
\toprule
Knob & Levels & Meaning \\
\midrule
$\gamma$ & $0,\,0.25,\,0.5,\,1,\,1.5$ & assignment strength (observed confounding) \\
$s$ & $0.5,\,1$ (and $0$) & hidden confounder $\to$ assignment \\
$h$ & $0,\,0.5,\,2$ (and $1$) & effect heterogeneity \\
\bottomrule
\end{tabular}
\end{table}

% -----------------------------------------------------------------------------------------
\subsection{Sweep results}
\label{app:semisynthetic:sweep-results}

All $30$ runs of Table~\ref{tab:semisynthetic:sweep} completed and were scored on all four
checkpoint kinds (\texttt{val\_loss}, \texttt{bestCI}, \texttt{bestCALIB}, \texttt{final\_epoch});
Table~\ref{tab:semisynthetic:checkpoint} pools every cell and replicate by kind. \texttt{val\_loss}
--- the real-data recommender's default --- has the highest pooled regret and PEHE of the four;
\texttt{bestCI} is lowest on both and is used for Table~\ref{tab:semisynthetic:sweep-results}.
This was a post-hoc choice using the reported hold-out's oracle metrics, while that hold-out
also supplied early stopping. The following tables are therefore exploratory historical
results, not independent validation. Protocol-v2 aggregation accepts a fixed checkpoint rule
(default \texttt{val\_loss}) and never selects it from the reported oracle regret.

\paragraph{Comparator interpretation.} The omitted Kaplan--Meier effect-error comparison
changes the conclusion: at these historical \texttt{bestCI} checkpoints DynaSurv has higher
mean pair PEHE in nine of ten cells. The reference values are $1.761$ versus $1.705$ months;
at hidden strength 1 they are $2.699$ versus $2.171$; at heterogeneity 0, $2.082$ versus
$1.192$. Only heterogeneity 2 favors DynaSurv ($2.643$ versus $2.773$). Reference policy
regret favors DynaSurv ($0.336$ versus $0.516$ months). The corrected aggregator reports
curve, effect, policy and factual metrics for every comparator on the same support. Ranking
benefit must not be interpreted as accurate individual treatment-effect estimation.

\begin{table}[ht]
\centering\small
\caption{Checkpoint-kind comparison, pooled over all 30 sweep runs. PEHE and $|$bias$|$
(months) are $n$-weighted means over arm pairs and cells; regret (months) and hit rate are
$n$-weighted means of the \texttt{line\_only} policy; factual $C$ is the IPCW concordance index.}
\label{tab:semisynthetic:checkpoint}
\begin{tabular}{@{}lrrrrr@{}}
\toprule
Checkpoint & PEHE & $|$bias$|$ & Regret & Hit rate & Factual $C$ \\
\midrule
\texttt{val\_loss} (real-data default) & 2.37 & 1.51 & 0.42 & 0.62 & 0.666 \\
\texttt{bestCI} & 1.95 & 1.14 & 0.36 & 0.64 & 0.675 \\
\texttt{bestCALIB} & 2.43 & 1.35 & 0.37 & 0.58 & 0.665 \\
\texttt{final\_epoch} & 2.23 & 1.16 & 0.40 & 0.52 & 0.663 \\
\bottomrule
\end{tabular}
\end{table}

\begin{table}[ht]
\centering\small
\caption{Sweep results at the \texttt{bestCI} checkpoint, mean $\pm$ std over 3 replicates.
PEHE and $|$bias$|$ (months) pool over lines and arm pairs; regret (months, oracle $=0$) and
hit rate use the \texttt{line\_only} policy; $h_r$ is the random policy's regret, an upper
bound on the task's headroom at that cell.}
\label{tab:semisynthetic:sweep-results}
\begin{tabular}{@{}llrrrrr@{}}
\toprule
Axis & Level & PEHE & $|$bias$|$ & Regret & Hit rate & $h_r$ \\
\midrule
$\gamma$ & 0        & $1.27\pm0.15$ & $0.34\pm0.15$ & $0.41\pm0.03$ & $0.57\pm0.05$ & 0.77 \\
$\gamma$ & 0.25      & $1.45\pm0.28$ & $0.70\pm0.40$ & $0.45\pm0.02$ & $0.56\pm0.02$ & 0.78 \\
$\gamma$ & 0.5       & $1.61\pm0.19$ & $0.85\pm0.23$ & $0.39\pm0.04$ & $0.59\pm0.04$ & 0.79 \\
$\gamma$ & 1 (ref.)  & $1.76\pm0.42$ & $0.87\pm0.33$ & $0.34\pm0.04$ & $0.64\pm0.04$ & 0.80 \\
$\gamma$ & 1.5       & $2.33\pm1.10$ & $1.49\pm1.03$ & $0.40\pm0.05$ & $0.62\pm0.02$ & 0.81 \\
$s$ & 0.5            & $2.25\pm0.32$ & $1.49\pm0.24$ & $0.44\pm0.04$ & $0.59\pm0.02$ & 0.80 \\
$s$ & 1               & $2.70\pm0.15$ & $2.14\pm0.16$ & $0.48\pm0.01$ & $0.57\pm0.01$ & 0.80 \\
$h$ & 0               & $2.08\pm0.75$ & $1.28\pm0.34$ & $0.03\pm0.03$ & $0.95\pm0.03$ & 0.46 \\
$h$ & 0.5             & $1.46\pm0.29$ & $0.94\pm0.45$ & $0.15\pm0.08$ & $0.67\pm0.12$ & 0.50 \\
$h$ & 2               & $2.64\pm0.53$ & $1.35\pm0.49$ & $0.52\pm0.07$ & $0.62\pm0.01$ & 1.46 \\
\bottomrule
\end{tabular}
\end{table}

\begin{table}[ht]
\centering\small
\caption{Overlap (effective sample size / $n$ per line, mean over 3 replicates), from the
synced manifests. Heterogeneity does not touch assignment, so it repeats the $\gamma=1$ row
exactly at every level.}
\label{tab:semisynthetic:overlap}
\begin{tabular}{@{}llrrrr@{}}
\toprule
Axis & Level & Line 1 & Line 2 & Line 3 & Line 4 \\
\midrule
$\gamma$ & 0        & 0.566 & 0.750 & 0.440 & 0.239 \\
$\gamma$ & 0.25      & 0.539 & 0.709 & 0.422 & 0.217 \\
$\gamma$ & 0.5       & 0.462 & 0.586 & 0.364 & 0.184 \\
$\gamma$ & 1 (ref.)  & 0.251 & 0.282 & 0.198 & 0.086 \\
$\gamma$ & 1.5       & 0.087 & 0.154 & 0.085 & 0.097 \\
$s$ & 0.5            & 0.197 & 0.275 & 0.140 & 0.061 \\
$s$ & 1               & 0.156 & 0.199 & 0.109 & 0.098 \\
$h$ & any level        & 0.251 & 0.282 & 0.198 & 0.086 \\
\bottomrule
\end{tabular}
\end{table}

\paragraph{Reading the three axes.} $\gamma$ (assignment on observed covariates) roughly
doubles PEHE and the average-effect bias from $\gamma=0$ to $\gamma=1.5$, but the model's own
regret is flat within noise (0.34--0.45) while the random-policy headroom $h_r$ grows only
slightly (0.77 to 0.81): the ranking of arms survives better than the size of the effect does.
This is consistent with overlap collapse: mean effective sample size / $n$ on line 4 falls from
$0.24$ at $\gamma=0$ to $0.10$ at $\gamma=1.5$ (Table~\ref{tab:semisynthetic:overlap}), so
precision, not direction, is what degrades. Three replicates cannot separate true bias from
finite-sample noise on this axis; the sweep gives a variance estimate, not a significance test.
$s$ (hidden confounder into assignment) is the sharper effect: bias more than doubles from the
$s=0$ reference ($0.87$) to $s=1$ ($2.14$) and regret rises by $40\%$ even though its overlap
loss is comparable to $\gamma$'s (line-4 ESS/$n$ falls to $0.06$--$0.10$) -- the pattern expected
of a confounder the model cannot see, not merely thinner support. $h$ (heterogeneity) mainly
changes the size of the task rather than the model's absolute accuracy: it does not change
assignment at all (identical overlap to the $\gamma=1$ reference row, Table~\ref{tab:semisynthetic:overlap}),
so any shift in the metrics below comes from the outcome model alone. At $h=0$ every line has
a single best arm, so both the model and a constant policy score near zero regret ($h_r=0.46$
is administrative-censoring headroom, not heterogeneity); as $h$ grows to $2$, $h_r$ triples to
$1.46$ and the model's regret grows with it (to $0.52$) while staying well under the random
policy's ceiling and the hit rate falls from $0.95$ to $0.62$ -- picking the right arm gets
harder as there are more genuinely different arms to pick between.

\paragraph{What this does not show.} The 30-run sweep uses one training configuration
(Section~\ref{app:semisynthetic:training}) and the model config predating this benchmark; it is
not re-tuned per cell. It has no Cox or DeepSurv baseline, so it cannot yet reproduce the
observed-vs-hidden-confounding bias comparison of Section~\ref{sec:counterfactual} with this
generator. Achieved censoring event rates matched their per-line targets to within $0.0007$ in
every one of the 30 synced manifests, confirming the calibration of
Section~\ref{app:semisynthetic:censoring} generalizes across the sweep.

% -----------------------------------------------------------------------------------------
\subsection{Verification of the generator}
\label{app:semisynthetic:verification}

The following were checked on the reference replicate.
\begin{itemize}
    \item Mean propensity per line equals the target to $10^{-6}$ for
    $\gamma\in\{0,0.5,1,2,4\}$ and $s\in\{0,2\}$; sampled arm counts pass a chi-square test;
    at $\gamma=s=0$ the propensity is constant within a line.
    \item A multinomial regression of the sampled arm on the drivers, pooled over lines with
    line indicators, recovers the assignment coefficients (correlation $0.976$, all signs
    correct). On line 1 alone it does not ($0.58$), because the previous-line drivers are
    constant there.
    \item Simulated Kaplan--Meier curves match the real ones (within about $0.04$ at 6, 12,
    24 and 36 months, medians within $2\%$); observed curves after censoring within about
    $0.02$; simulated event rates equal the real rates.
    \item Closed-form RMST agrees with Monte Carlo (bias $-0.002$ months, standardised
    differences of standard deviation $1.01$); at $h=0$, $\tau_a=0$ all arms have equal $\eta$.
    \item Expansion: the unmasked rows per line equal the real patients per line; history rows
    equal the real records; no hidden variable and no propensity in the model-facing files.
    \item Reproducibility: an identical rerun is bit-identical; changing $\gamma$ keeps $u$
    and changes $13\%$ of the arms; a new replicate changes $u$.
    \item No patient appears on both sides of a temporal, random or cross-validation split.
\end{itemize}

% -----------------------------------------------------------------------------------------
\subsection{Limitations}
\label{app:semisynthetic:limitations}

\begin{enumerate}
    \item Coefficients are hand-set; results describe the difficulty of this generator, not
    real effect sizes.
    \item One Weibull shape per line, a shared coupling across arms, non-informative censoring,
    no outcome feedback across lines, and no treatment change within a line.
    \item Only $16$ covariates affect the outcome; the model sees many more real columns, which
    are non-causal here by construction.
    \item Arms are thin on late lines (\arm{ET alone} and \arm{ET+anti-CDK} on line 4): their
    counterfactuals are extrapolation, and unsupported arms are excluded from the metrics.
    \item Calendar time is an instrument and the hold-out enters late: an intentional shift
    that the reported errors include.
    \item The covariate standardiser is fitted on all samples, hold-out included (a property of
    the existing data module).
    \item Training settings were tuned on real data and chosen by hand here; the benchmark
    scores the model \emph{as trained}, not the best model it could be.
\end{enumerate}

% Needs only amsmath and booktabs, both already loaded there. Self-contained: no \todo,
% no external figures. Code references are to src/CausalSurv/semisynthetic/.

\providecommand{\drv}[1]{\texttt{#1}}
\providecommand{\arm}[1]{\textsc{#1}}
\providecommand{\Ind}{\mathbf{1}}

\section{Semi-Synthetic Data Generation}
\label{app:semisynthetic}

This appendix records every design choice of the semi-synthetic benchmark used for
counterfactual validation (Section~\ref{sec:counterfactual}): what was decided, and why.
The benchmark keeps the real covariates and treatment history of the cohort, and
\emph{simulates} the treatment of the line being scored, the potential survival time under
every arm, and the censoring. The true survival curve, restricted mean survival time (RMST)
and best arm of every patient under every arm are then known in closed form.

The numbers quoted below are for the reference configuration (assignment strength
$\gamma=1$, hidden-confounder strength $s=0$, heterogeneity $h=1$). All coefficients are
set by hand, not estimated: they fix the \emph{difficulty} of the benchmark, not its realism.
DynaSurv's own results on the sweep (Section~\ref{app:semisynthetic:sweep}) are reported
against the oracle and naive Kaplan--Meier references only: this rebuild does not yet include
the Cox/DeepSurv baselines used for the observed-vs-hidden-confounding argument in
Section~\ref{sec:counterfactual}.

% -----------------------------------------------------------------------------------------
\subsection{Cohort and arms}
\label{app:semisynthetic:cohort}

\paragraph{Cohort.} HR$^+$/HER2$^-$ patients whose first treatment line starts in 2018 or
later, lines 1--4, built exactly as the training pipeline builds its cohort (patient-level
entry-year filter, inner join with the static file), so that simulated rows line up one to
one with the rows the model trains on. This gives $6{,}191$ patients and $13{,}116$
(patient, line) records, $6{,}191 / 3{,}401 / 2{,}169 / 1{,}355$ on lines 1 to 4. The year
restriction is the identifiability window of the real pipeline: every arm was available
throughout. One exact duplicate (patient, line) record in the source file is removed, because
after expansion it would become a sample of its own with a false ``previous line''.
Administrative censoring is measured from the last observed line start, 2024-02-21.

\paragraph{Arms.} Four arms: \arm{ET alone}, \arm{ET+anti-CDK wo CT}, \arm{MonoCT std alone},
\arm{PolyCT alone}. \arm{CT+anti-angio} is not simulated (21, 7, 2 and 2 records per line in the
restricted cohort: no overlap). It and all other categories (\arm{CT+anti-HER2}, \arm{CT+IT},
\arm{CT+TT}, \arm{ET+TT}, \arm{No treatment}, \arm{Other}) remain in the data as
\emph{history} only and are excluded from the set of recommendable arms. Arm indices follow
the alphabetical order of the arm names throughout.

\paragraph{Real history, simulated present.} Covariates, past treatments and line structure
(previous-line duration, gap, line number) are real. Only the arm and outcome of the line
being scored are simulated, and the simulated outcome does \emph{not} feed back into later
lines. This keeps the covariate distribution real and makes each line a separate problem
given its history. The price is that the benchmark does not test how a model handles
treatment--outcome feedback across lines.

% -----------------------------------------------------------------------------------------
\subsection{Prefix expansion}
\label{app:semisynthetic:expansion}

A patient with $L$ lines becomes $L$ samples. Sample $j\in\{1,\dots,L\}$ contains the real
records of lines $1,\dots,j-1$ and record $j$ with the simulated arm and simulated outcome;
its identifier is $10\cdot\text{id}+j$, and all samples of a patient share the original id
(the unit of any train/validation grouping) and the entry year. The reference cohort yields
$13{,}116$ samples in $24{,}920$ rows.

The reason is that the sequence model steps through lines in order, and its hidden state at
line $j$ does not depend on the arm at line $j$. Scoring only the last line of each sample
therefore requires no change to the model or its losses; the alternative, simulating all
lines of a patient jointly, would require outcome feedback and a joint model. The data
module masks every row except the sample's last line; history rows carry placeholder
outcomes ($0$).

Columns describing the real current line and its outcome (drug flags other than the treatment
category, all outcome columns other than the two survival targets, line end date, death date)
are removed, so that no real value sits next to a simulated one. Static features do not reach
the encoder, so age and menopausal status are copied into dynamic columns; otherwise they
would act as unplanned hidden confounders.

% -----------------------------------------------------------------------------------------
\subsection{Drivers and causal roles}
\label{app:semisynthetic:drivers}

Each sample has a vector $z\in\mathbb{R}^{16}$ of drivers computed from its real record.
Continuous drivers are standardised over the cohort; the others are $0/1$ flags.
Table~\ref{tab:semisynthetic:drivers} gives their definitions and roles. The role of each
driver is enforced when the configuration is loaded: a coefficient on an instrument in the
outcome equation, or on a prognostic factor in the assignment equation, is rejected.

\begin{table}[ht]
\centering\small
\caption{Drivers of the data-generating process. ``Both'' means the driver enters assignment
and outcome (confounder).}
\label{tab:semisynthetic:drivers}
\begin{tabular}{@{}lll@{}}
\toprule
Driver & Definition & Role \\
\midrule
\drv{liver} & liver metastases $>0$ & both \\
\drv{visceral} & pulmonary or pleural metastases $>0$ & both \\
\drv{bone\_only} & bone metastases, none of liver, visceral, brain & both \\
\drv{mpps} & imputed ECOG/WHO performance status, 0--4 (standardised) & both \\
\drv{log\_prev\_line} & $\log(1+\text{previous line duration})$ (standardised) & both \\
\drv{prev\_cdk}, \drv{prev\_et}, \drv{prev\_ct} & class of the \emph{real} previous arm (one-hot; all zero at line 1) & both \\
\drv{cum\_new\_sites} & cumulative new metastatic sites (standardised) & both \\
\drv{age}, \drv{menopause} & age at selection (standardised), menopausal status & both \\
\drv{calendar} & months since 2018-01-01 (standardised) & assignment only \\
\drv{brain} & brain metastases $>0$ & outcome only \\
\drv{lobular} & invasive lobular histology & outcome only \\
\drv{brca} & \textit{BRCA1} or \textit{BRCA2} & outcome only \\
\drv{old\_site\_progression} & progression of an old site $>0$ & outcome only \\
\bottomrule
\end{tabular}
\end{table}

The model still receives all real covariate columns ($73$, plus the calendar covariate); only
these $16$ matter for the simulated data, and the remainder are nuisance columns it must
learn to ignore. A buffer-time covariate is deliberately not a driver.

\paragraph{Hidden confounder.} A latent variable
$u=\rho\,\tilde z_{\text{liver}}+\sqrt{1-\rho^2}\,\varepsilon$ with $\varepsilon\sim\mathcal{N}(0,1)$,
$\tilde z_{\text{liver}}$ the standardised liver indicator and $\rho=0.3$, enters both
assignment and outcome and is \emph{never} written to the model input. Correlating it with a
real covariate keeps it from being independent noise. Its outcome coefficient is fixed and
only its coupling to assignment is swept, so that strength $0$ is an unconfounded reference
carrying the same outcome heterogeneity: any error growth above it is confounding bias.

% -----------------------------------------------------------------------------------------
\subsection{Treatment assignment}
\label{app:semisynthetic:assignment}

For a sample on line $\ell$ with drivers $z$ and hidden $u$, the arm $A$ is drawn with
\begin{equation}
    \pi_a(z,u)=\mathbb{P}(A=a\mid z,u)
    =\frac{\exp\eta^{\text{A}}_a}{\sum_{b}\exp\eta^{\text{A}}_b},
    \qquad
    \eta^{\text{A}}_a=\alpha_{\ell a}+\gamma\,B_a^{\top}z+s\,c_a\,u,
    \label{eq:semisynthetic:assignment}
\end{equation}
where $B_a$ are hand-set coefficients (Table~\ref{tab:semisynthetic:coefficients}, shared
across lines), $c_a$ the hidden loadings, $\gamma$ the assignment strength, and $s$ the hidden
strength. Coefficients are chosen for clinically plausible signs (visceral involvement pushes
towards combination chemotherapy, a long previous line towards endocrine-based therapy, a CDK
inhibitor is rarely repeated) and moderate size.

\paragraph{Calibrated intercepts.} The intercepts $\alpha_{\ell a}$ are not set by hand. For
every line they are fitted, by iterative proportional fitting
\begin{equation}
    \alpha_{\ell a}\leftarrow\alpha_{\ell a}+\log\frac{q_{\ell a}}{\bar\pi_{\ell a}},
    \qquad
    \bar\pi_{\ell a}=\frac{1}{n_\ell}\sum_{i\in\ell}\pi_a(z_i,u_i),
\end{equation}
until $\max_a|\bar\pi_{\ell a}-q_{\ell a}|<10^{-6}$, where $q_{\ell\cdot}$ is the real
distribution of the four arms on line $\ell$. The marginal arm mix therefore equals the real
mix at every $\gamma$ and $s$: a sweep changes \emph{who} receives which arm, never \emph{how
many}. The average runs over every sample of the line (each is given a simulated arm) while
the target is the real mix among the four arms only.
The targets $q_\ell$ (order: ET, ET+CDK, MonoCT, PolyCT) are
$(.183,.628,.090,.100)$, $(.148,.412,.330,.109)$, $(.049,.107,.611,.233)$ and
$(.032,.031,.580,.357)$ for lines 1 to 4.

\paragraph{Sampling.} The arm is drawn by inverting the cumulative distribution with
\emph{one uniform per sample}, so that draws coincide wherever propensities coincide. The true
propensities are stored.

\paragraph{Range of the sweep.} The effective sample size of inverse-propensity weights of the
assigned arm, divided by $n$, is $0.51$, $0.40$, $0.22$, $0.01$ and $0.00$ at $\gamma=0$, $0.5$,
$1$, $2$ and $4$ respectively (no hidden confounder), and $0.07$ at $s=2$ even with $\gamma=0$. Beyond $\gamma\approx1.5$ or
$s\approx1$ the benchmark measures extrapolation rather than confounding, so the sweep is
restricted to $\gamma\le1.5$ and $s\le1$.

\begin{table}[ht]
\centering\small
\caption{Hand-set coefficients (blank $=0$). Left: assignment coefficients $B_a$ and hidden
loadings $c_a$. Right: outcome prognostic term $f$ and arm interactions $g_a$. Continuous
drivers are per standard deviation, flags are on \emph{vs} off; outcome coefficients are on the
log-hazard scale (positive $=$ worse). The row ``Hidden'' gives the loadings $c_a$ of $u$ on the
arms' assignment logits and, in the $f$ column, its outcome coefficient $\beta_u$.}
\label{tab:semisynthetic:coefficients}
\setlength{\tabcolsep}{4pt}
\begin{tabular}{@{}l rrrr c rrrrr@{}}
\toprule
 & \multicolumn{4}{c}{Assignment $B_a$} & & \multicolumn{5}{c}{Outcome} \\
\cmidrule(lr){2-5}\cmidrule(l){7-11}
Driver & ET & CDK & Mono & Poly & & $f$ & $g_{\text{ET}}$ & $g_{\text{CDK}}$ & $g_{\text{Mono}}$ & $g_{\text{Poly}}$ \\
\midrule
\drv{liver}          &      & $-0.3$ & $0.5$  & $0.9$  & & $0.45$ & $0.30$ & $0.25$ & $-0.15$ & $-0.35$ \\
\drv{visceral}       &      & $-0.2$ & $0.3$  & $0.5$  & & $0.25$ &        &        &         & $-0.25$ \\
\drv{bone\_only}     &      & $0.5$  & $-0.4$ & $-0.6$ & & $-0.30$ & $-0.15$ & $-0.30$ &      &         \\
\drv{mpps}           &      & $-0.3$ & $0.3$  & $-0.2$ & & $0.30$ &        &        & $-0.10$ & $0.15$  \\
\drv{log\_prev\_line} &     & $0.4$  & $-0.4$ & $-0.5$ & & $-0.25$ &       & $-0.20$ &        &         \\
\drv{prev\_cdk}      & $-0.3$ & $-1.0$ & $0.4$ & $0.3$  & & $0.20$ &        & $0.40$ &         &         \\
\drv{prev\_et}       & $-0.5$ &      &        &        & &        & $0.30$ &        &         &         \\
\drv{prev\_ct}       &      &        & $0.3$  &        & &        &        &        &         &         \\
\drv{cum\_new\_sites} &     & $-0.1$ & $0.2$  & $0.3$  & & $0.15$ &        &        &         &         \\
\drv{age}            & $0.4$ &       & $0.3$  & $-0.4$ & & $0.10$ &        &        & $-0.10$ & $0.15$  \\
\drv{menopause}      &      & $0.3$  &        & $-0.3$ & & $0.05$ &        &        &         &         \\
\drv{calendar}       &      & $0.3$  & $-0.1$ & $-0.2$ & &        &        &        &         &         \\
\drv{brain}          &      &        &        &        & & $0.40$ &        &        &         &         \\
\drv{lobular}        &      &        &        &        & & $0.10$ &        &        &         &         \\
\drv{brca}           &      &        &        &        & & $0.15$ &        &        &         &         \\
\drv{old\_site\_prog.} &    &        &        &        & & $0.15$ &        &        &         &         \\
\midrule
Hidden ($c_a$; $\beta_u$ under $f$) & 0 & $-0.5$ & $0.5$ & $1.0$ & & $0.4$ & & & & \\
Arm effect $\tau_a$  &      &        &        &        & & & $0$ & $-0.20$ & $0.10$ & $0$ \\
\bottomrule
\end{tabular}
\end{table}

% -----------------------------------------------------------------------------------------
\subsection{Potential outcomes}
\label{app:semisynthetic:outcome}

Time from the start of the line to death under arm $a$ follows a Weibull model with
line-specific shape $k_\ell$ and scale $\lambda_\ell$,
\begin{equation}
    S_a(t\mid z,u)=\exp\!\Big(-\big(t/\lambda_\ell\big)^{k_\ell}e^{\eta_a}\Big),
    \qquad
    \eta_a=f(z)+\tau_a+h\,g_a(z)+\beta_u\,u,
    \label{eq:semisynthetic:outcome}
\end{equation}
that is, $T_a\sim\text{Weibull}(k_\ell,\lambda_\ell e^{-\eta_a/k_\ell})$; $\eta_a>0$ is a higher
hazard. $f$ is the prognostic term, $\tau_a$ the arm effect, $g_a$ the interaction that makes
the best arm differ across patients (who benefits from what), $h$ the heterogeneity, and
$\beta_u=0.4$ the fixed outcome coefficient of the hidden confounder. The functions $f$ and
$g_a$ are linear in $z$ (Table~\ref{tab:semisynthetic:coefficients}) and \emph{centred within
each line}, which makes $(k_\ell,\lambda_\ell)$ the baseline of the average patient.

\paragraph{Fitting the baseline.} For each line, $(k_\ell,\lambda_\ell)$ minimise the squared
difference between the population survival implied by the simulation for the assigned arms,
$\bar S_\ell(t)=n_\ell^{-1}\sum_{i\in\ell}\exp\!\big(-(t/\lambda)^{k}e^{\eta_{i,A_i}}\big)$, and the
real per-line Kaplan--Meier curve, over the observed jump times at which at least $5\%$ of the
line is still at risk. The result is $k_\ell\approx1.27$--$1.33$ and
$\lambda_\ell\approx53/27/22/18$ months; the fitted medians are $48.5/22.6/16.3/12.7$ months
against $47.7/22.8/16.5/12.4$ in the real data.

\paragraph{Closed-form truth.} Survival is given by \eqref{eq:semisynthetic:outcome}. With
$\sigma_a=\lambda_\ell e^{-\eta_a/k_\ell}$ the RMST to horizon $\tau$ is
\begin{equation}
    \text{RMST}_a(\tau)=\int_0^{\tau}S_a(t)\,dt
    =\sigma_a\,\Gamma\!\big(1+\tfrac1{k_\ell}\big)\,P\!\big(\tfrac1{k_\ell},(\tau/\sigma_a)^{k_\ell}\big),
\end{equation}
with $P$ the regularised lower incomplete gamma function; the individual treatment effect of
arm $a$ against $b$ is $\text{RMST}_a-\text{RMST}_b$ and the best arm maximises RMST over the
recommendable arms.

\paragraph{Coupling of the arms.} A single uniform per sample is shared by all arms, so
$T_a=\sigma_a(-\log V)^{1/k_\ell}$ with the same $V\sim\mathcal{U}(0,1)$ for every $a$. The
arms then differ only through $\eta_a$, and the individual effect is not swamped by independent
noise. The marginals are unaffected; the joint distribution of potential outcomes is one
particular choice, and another coupling would give the same marginals with different
individual-level effects.

\paragraph{Effect sizes.} Effects are small on the RMST scale: the average true pairwise
differences at the reference configuration lie between $-0.44$ and $+0.42$ months, and the
\emph{headroom} of a policy (value of the oracle minus value of the best single arm per line)
is $0.00, 0.11, 0.47, 1.25, 1.71$ months for $h=0,0.5,1,2,3$. The arm effect of
\arm{ET+anti-CDK} was first $-0.35$, which made it the best arm for $63$--$74\%$ of patients
(by line) so that ``always CDK'' scored highly; at $-0.20$ its share on line 2 drops from
$63\%$ to $48\%$ and the best-arm metric discriminates.

% -----------------------------------------------------------------------------------------
\subsection{Censoring}
\label{app:semisynthetic:censoring}

The observed time and event indicator are $\min(T,C)$ and $\Ind\{T\le C\}$ with
\begin{equation}
    C=\min\big(C^{\text{adm}},\,C^{\text{drop}}\big),\qquad
    C^{\text{adm}}=\frac{t_{\text{cutoff}}-t^{\text{start}}}{30.44},\qquad
    C^{\text{drop}}\sim\text{Exp}(r_\ell),
\end{equation}
in months, with $T$ the latent time under the simulated arm. Administrative censoring is a
fact of the real cohort and is kept sample by sample. Dropout is added only where the cutoff
alone leaves the simulated event rate above the real one; $r_\ell$ is found by bisection on
$\log r$ (60 steps) so that the event rate equals the real rate of the restricted cohort,
using fixed unit-exponential draws so that the event rate is a deterministic decreasing
function of $r$. The targets are the event rates of the restricted cohort,
$0.442/0.547/0.578/0.596$ (not those of the full 2008-onwards cohort, $0.73$--$0.85$). With
the cutoff alone the rates would be $0.465/0.587/0.639/0.679$, so dropout is active on every
line, at $0.0024/0.0041/0.0073/0.0103$ per month.

Dropout draws are independent conditional on line and the fitted generator parameters.
Administrative follow-up depends on calendar date, which also affects assignment and is
associated with patient characteristics. Marginal independent censoring is therefore not
guaranteed, and training-cohort censoring weights need not transfer to the later test cohort.
Protocol v2 uses known survival probabilities and uncensored latent outcomes for primary
simulation scoring instead of making that transfer assumption.

% -----------------------------------------------------------------------------------------
\subsection{Randomness and artefacts}
\label{app:semisynthetic:seeds}

Four independent random streams (hidden confounder, assignment, outcome, censoring) are
spawned from one seed sequence keyed by (data seed, replicate). A sweep over $\gamma$, $s$ or
$h$ therefore reuses the same $u$, the same assignment uniforms, the same latent uniforms and
the same dropout draws: levels differ \emph{only} by the swept knob, and repeating a
configuration reproduces the data exactly.

Each replicate writes (i) the expanded dynamic and re-keyed static tables under the names of
the real data files, so that the directory can be used in place of the real data; (ii) a truth
table with, per sample, the assigned arm, hidden $u$, Weibull parameters, latent time under
every arm, both censoring times, observed time and event, and propensity and $\eta$ for every
arm; and (iii) a manifest with every parameter, the calibrated intercepts, the arm counts per
line, the fitted Weibull parameters and medians, the dropout rates, and the achieved and
target event rates. The hidden $u$ appears only in the truth table.

% -----------------------------------------------------------------------------------------
\subsection{Training and data-loading choices}
\label{app:semisynthetic:training}

\begin{itemize}
    \item \textbf{Mask.} Only the last line of each sample is unmasked, and the counts that
    decide which arms are recommendable are computed on those rows.
    \item \textbf{Grouped splits.} Every split other than the temporal one (random hold-out,
    cross-validation folds, early-stopping split) groups by original patient, because the
    samples of a patient share their history and would otherwise leak across the boundary. A
    temporal split already keeps them together (shared entry year).
    \item \textbf{Protocol-v2 validation.} Twenty percent of the pre-2021 original patients
    form a separate development validation set (fixed split seed 0); the remainder train
    the model, scaler and support models. The survival grid is fitted from training endpoint
    outcomes only. The temporal test set never supplies early stopping or hyperparameter
    selection. Historical results below predate this correction and require rerunning.
    \item \textbf{Temporal hold-out.} Patients entering in 2021 or later ($26\%$ of samples),
    a deliberate policy-era shift; calendar time is an instrument of the assignment, so the
    model receives it.
    \item \textbf{Support thresholds.} Arms need at least $30$ samples, $10$ events and $30$
    samples followed to the horizon at a line to be recommendable, against $200/30/100$ on
    real data, because the four arms thin out on later lines (on line 4, \arm{ET alone} has
    $25$ samples and is excluded there).
    \item \textbf{Optimisation.} Early stopping on the validation loss (patience $20$, at most
    $150$ epochs). The architecture and learning-rate settings are those tuned on the real
    data; they were not re-tuned for the synthetic data.
\end{itemize}

% -----------------------------------------------------------------------------------------
\subsection{Evaluation against the truth}
\label{app:semisynthetic:evaluation}

All quantities are computed on the temporal hold-out, at each sample's own line $\ell$ and its
horizon $\tau_\ell$ ($24, 18, 12, 12$ months), on a $0.1$-month grid, and only for arms that
are recommendable at that line. Write $\hat S_a$ for a predicted curve and $S_a$ for the truth.
\begin{itemize}
    \item \textbf{Curve error.} $\big(\tau_\ell^{-1}\!\int_0^{\tau_\ell}(\hat S_a-S_a)^2dt\big)^{1/2}$
    averaged in squares over samples, and the error and bias of $\widehat{\text{RMST}}_a-\text{RMST}_a$,
    reported separately for the assigned (factual) arm and for the other (counterfactual) arms.
    \item \textbf{Effect error.} For each arm pair $(a,b)$,
    $\text{PEHE}=\big(\mathbb{E}[(\hat\Delta_{ab}-\Delta_{ab})^2]\big)^{1/2}$ with
    $\Delta_{ab}=\text{RMST}_a-\text{RMST}_b$, the bias of the average effect
    $\mathbb{E}[\hat\Delta_{ab}]-\mathbb{E}[\Delta_{ab}]$, and the share of samples with the
    correct sign. Hidden confounding can affect both bias and PEHE.
    \item \textbf{Policy.} A policy picks the arm with the highest predicted RMST among
    supported arms (abstaining to the assigned arm where none is supported). Value is the
    true RMST of the chosen arm; regret is the oracle value minus it; the hit rate is the share
    of samples where the chosen arm is the true best arm. Reference policies: the assigned
    arm, the best single arm per line (an upper bound for constant policies), and uniform
    random.
    \item \textbf{Factual accuracy (protocol v2).} At the exact requested horizon, write
    $p=S_A(\tau_\ell)$ and $q=\hat S_A(\tau_\ell)$. The primary Brier score averages
    $p(1-q)^2+(1-p)q^2$. We also report Brier scores from uncensored latent outcomes and
    concordance of latent event times with horizon-specific predicted risk. These are
    simulation diagnostics, not estimators available on real censored data. The historical
    IPCW results below used training-cohort censoring weights and must be replaced.
    \item \textbf{References.} The \emph{oracle} (true curves; must score zero error and zero
    regret) and the \emph{naive Kaplan--Meier} (per line and arm, no covariates; the bar a
    covariate model must clear).
\end{itemize}

% -----------------------------------------------------------------------------------------
\subsection{Sweep design}
\label{app:semisynthetic:sweep}

One knob is varied at a time, the others staying at the reference values.
Table~\ref{tab:semisynthetic:sweep} lists the cells; each is replicated three times (data seed
and training seed equal to the replicate index), giving $30$ runs. The hidden-confounder axis
is the one that separates methods that adjust for observed covariates only from methods that
would be unbiased under hidden confounding; the assignment axis is expected to be flat for
outcome-regression estimators, since observed confounding adds variance rather than bias.

\begin{table}[ht]
\centering\small
\caption{Sweep cells. The reference cell is $\gamma=1$, $s=0$, $h=1$; $s=0$ and $h=1$ are
therefore not repeated.}
\label{tab:semisynthetic:sweep}
\begin{tabular}{@{}lll@{}}
\toprule
Knob & Levels & Meaning \\
\midrule
$\gamma$ & $0,\,0.25,\,0.5,\,1,\,1.5$ & assignment strength (observed confounding) \\
$s$ & $0.5,\,1$ (and $0$) & hidden confounder $\to$ assignment \\
$h$ & $0,\,0.5,\,2$ (and $1$) & effect heterogeneity \\
\bottomrule
\end{tabular}
\end{table}

% -----------------------------------------------------------------------------------------
\subsection{Sweep results}
\label{app:semisynthetic:sweep-results}

All $30$ runs of Table~\ref{tab:semisynthetic:sweep} completed and were scored on all four
checkpoint kinds (\texttt{val\_loss}, \texttt{bestCI}, \texttt{bestCALIB}, \texttt{final\_epoch});
Table~\ref{tab:semisynthetic:checkpoint} pools every cell and replicate by kind. \texttt{val\_loss}
--- the real-data recommender's default --- has the highest pooled regret and PEHE of the four;
\texttt{bestCI} is lowest on both and is used for Table~\ref{tab:semisynthetic:sweep-results}.
This was a post-hoc choice using the reported hold-out's oracle metrics, while that hold-out
also supplied early stopping. The following tables are therefore exploratory historical
results, not independent validation. Protocol-v2 aggregation accepts a fixed checkpoint rule
(default \texttt{val\_loss}) and never selects it from the reported oracle regret.

\paragraph{Comparator interpretation.} The omitted Kaplan--Meier effect-error comparison
changes the conclusion: at these historical \texttt{bestCI} checkpoints DynaSurv has higher
mean pair PEHE in nine of ten cells. The reference values are $1.761$ versus $1.705$ months;
at hidden strength 1 they are $2.699$ versus $2.171$; at heterogeneity 0, $2.082$ versus
$1.192$. Only heterogeneity 2 favors DynaSurv ($2.643$ versus $2.773$). Reference policy
regret favors DynaSurv ($0.336$ versus $0.516$ months). The corrected aggregator reports
curve, effect, policy and factual metrics for every comparator on the same support. Ranking
benefit must not be interpreted as accurate individual treatment-effect estimation.

\begin{table}[ht]
\centering\small
\caption{Checkpoint-kind comparison, pooled over all 30 sweep runs. PEHE and $|$bias$|$
(months) are $n$-weighted means over arm pairs and cells; regret (months) and hit rate are
$n$-weighted means of the \texttt{line\_only} policy; factual $C$ is the IPCW concordance index.}
\label{tab:semisynthetic:checkpoint}
\begin{tabular}{@{}lrrrrr@{}}
\toprule
Checkpoint & PEHE & $|$bias$|$ & Regret & Hit rate & Factual $C$ \\
\midrule
\texttt{val\_loss} (real-data default) & 2.37 & 1.51 & 0.42 & 0.62 & 0.666 \\
\texttt{bestCI} & 1.95 & 1.14 & 0.36 & 0.64 & 0.675 \\
\texttt{bestCALIB} & 2.43 & 1.35 & 0.37 & 0.58 & 0.665 \\
\texttt{final\_epoch} & 2.23 & 1.16 & 0.40 & 0.52 & 0.663 \\
\bottomrule
\end{tabular}
\end{table}

\begin{table}[ht]
\centering\small
\caption{Sweep results at the \texttt{bestCI} checkpoint, mean $\pm$ std over 3 replicates.
PEHE and $|$bias$|$ (months) pool over lines and arm pairs; regret (months, oracle $=0$) and
hit rate use the \texttt{line\_only} policy; $h_r$ is the random policy's regret, an upper
bound on the task's headroom at that cell.}
\label{tab:semisynthetic:sweep-results}
\begin{tabular}{@{}llrrrrr@{}}
\toprule
Axis & Level & PEHE & $|$bias$|$ & Regret & Hit rate & $h_r$ \\
\midrule
$\gamma$ & 0        & $1.27\pm0.15$ & $0.34\pm0.15$ & $0.41\pm0.03$ & $0.57\pm0.05$ & 0.77 \\
$\gamma$ & 0.25      & $1.45\pm0.28$ & $0.70\pm0.40$ & $0.45\pm0.02$ & $0.56\pm0.02$ & 0.78 \\
$\gamma$ & 0.5       & $1.61\pm0.19$ & $0.85\pm0.23$ & $0.39\pm0.04$ & $0.59\pm0.04$ & 0.79 \\
$\gamma$ & 1 (ref.)  & $1.76\pm0.42$ & $0.87\pm0.33$ & $0.34\pm0.04$ & $0.64\pm0.04$ & 0.80 \\
$\gamma$ & 1.5       & $2.33\pm1.10$ & $1.49\pm1.03$ & $0.40\pm0.05$ & $0.62\pm0.02$ & 0.81 \\
$s$ & 0.5            & $2.25\pm0.32$ & $1.49\pm0.24$ & $0.44\pm0.04$ & $0.59\pm0.02$ & 0.80 \\
$s$ & 1               & $2.70\pm0.15$ & $2.14\pm0.16$ & $0.48\pm0.01$ & $0.57\pm0.01$ & 0.80 \\
$h$ & 0               & $2.08\pm0.75$ & $1.28\pm0.34$ & $0.03\pm0.03$ & $0.95\pm0.03$ & 0.46 \\
$h$ & 0.5             & $1.46\pm0.29$ & $0.94\pm0.45$ & $0.15\pm0.08$ & $0.67\pm0.12$ & 0.50 \\
$h$ & 2               & $2.64\pm0.53$ & $1.35\pm0.49$ & $0.52\pm0.07$ & $0.62\pm0.01$ & 1.46 \\
\bottomrule
\end{tabular}
\end{table}

\begin{table}[ht]
\centering\small
\caption{Overlap (effective sample size / $n$ per line, mean over 3 replicates), from the
synced manifests. Heterogeneity does not touch assignment, so it repeats the $\gamma=1$ row
exactly at every level.}
\label{tab:semisynthetic:overlap}
\begin{tabular}{@{}llrrrr@{}}
\toprule
Axis & Level & Line 1 & Line 2 & Line 3 & Line 4 \\
\midrule
$\gamma$ & 0        & 0.566 & 0.750 & 0.440 & 0.239 \\
$\gamma$ & 0.25      & 0.539 & 0.709 & 0.422 & 0.217 \\
$\gamma$ & 0.5       & 0.462 & 0.586 & 0.364 & 0.184 \\
$\gamma$ & 1 (ref.)  & 0.251 & 0.282 & 0.198 & 0.086 \\
$\gamma$ & 1.5       & 0.087 & 0.154 & 0.085 & 0.097 \\
$s$ & 0.5            & 0.197 & 0.275 & 0.140 & 0.061 \\
$s$ & 1               & 0.156 & 0.199 & 0.109 & 0.098 \\
$h$ & any level        & 0.251 & 0.282 & 0.198 & 0.086 \\
\bottomrule
\end{tabular}
\end{table}

\paragraph{Reading the three axes.} $\gamma$ (assignment on observed covariates) roughly
doubles PEHE and the average-effect bias from $\gamma=0$ to $\gamma=1.5$, but the model's own
regret is flat within noise (0.34--0.45) while the random-policy headroom $h_r$ grows only
slightly (0.77 to 0.81): the ranking of arms survives better than the size of the effect does.
This is consistent with overlap collapse: mean effective sample size / $n$ on line 4 falls from
$0.24$ at $\gamma=0$ to $0.10$ at $\gamma=1.5$ (Table~\ref{tab:semisynthetic:overlap}), so
precision, not direction, is what degrades. Three replicates cannot separate true bias from
finite-sample noise on this axis; the sweep gives a variance estimate, not a significance test.
$s$ (hidden confounder into assignment) is the sharper effect: bias more than doubles from the
$s=0$ reference ($0.87$) to $s=1$ ($2.14$) and regret rises by $40\%$ even though its overlap
loss is comparable to $\gamma$'s (line-4 ESS/$n$ falls to $0.06$--$0.10$) -- the pattern expected
of a confounder the model cannot see, not merely thinner support. $h$ (heterogeneity) mainly
changes the size of the task rather than the model's absolute accuracy: it does not change
assignment at all (identical overlap to the $\gamma=1$ reference row, Table~\ref{tab:semisynthetic:overlap}),
so any shift in the metrics below comes from the outcome model alone. At $h=0$ every line has
a single best arm, so both the model and a constant policy score near zero regret ($h_r=0.46$
is administrative-censoring headroom, not heterogeneity); as $h$ grows to $2$, $h_r$ triples to
$1.46$ and the model's regret grows with it (to $0.52$) while staying well under the random
policy's ceiling and the hit rate falls from $0.95$ to $0.62$ -- picking the right arm gets
harder as there are more genuinely different arms to pick between.

\paragraph{What this does not show.} The 30-run sweep uses one training configuration
(Section~\ref{app:semisynthetic:training}) and the model config predating this benchmark; it is
not re-tuned per cell. It has no Cox or DeepSurv baseline, so it cannot yet reproduce the
observed-vs-hidden-confounding bias comparison of Section~\ref{sec:counterfactual} with this
generator. Achieved censoring event rates matched their per-line targets to within $0.0007$ in
every one of the 30 synced manifests, confirming the calibration of
Section~\ref{app:semisynthetic:censoring} generalizes across the sweep.

% -----------------------------------------------------------------------------------------
\subsection{Verification of the generator}
\label{app:semisynthetic:verification}

The following were checked on the reference replicate.
\begin{itemize}
    \item Mean propensity per line equals the target to $10^{-6}$ for
    $\gamma\in\{0,0.5,1,2,4\}$ and $s\in\{0,2\}$; sampled arm counts pass a chi-square test;
    at $\gamma=s=0$ the propensity is constant within a line.
    \item A multinomial regression of the sampled arm on the drivers, pooled over lines with
    line indicators, recovers the assignment coefficients (correlation $0.976$, all signs
    correct). On line 1 alone it does not ($0.58$), because the previous-line drivers are
    constant there.
    \item Simulated Kaplan--Meier curves match the real ones (within about $0.04$ at 6, 12,
    24 and 36 months, medians within $2\%$); observed curves after censoring within about
    $0.02$; simulated event rates equal the real rates.
    \item Closed-form RMST agrees with Monte Carlo (bias $-0.002$ months, standardised
    differences of standard deviation $1.01$); at $h=0$, $\tau_a=0$ all arms have equal $\eta$.
    \item Expansion: the unmasked rows per line equal the real patients per line; history rows
    equal the real records; no hidden variable and no propensity in the model-facing files.
    \item Reproducibility: an identical rerun is bit-identical; changing $\gamma$ keeps $u$
    and changes $13\%$ of the arms; a new replicate changes $u$.
    \item No patient appears on both sides of a temporal, random or cross-validation split.
\end{itemize}

% -----------------------------------------------------------------------------------------
\subsection{Limitations}
\label{app:semisynthetic:limitations}

\begin{enumerate}
    \item Coefficients are hand-set; results describe the difficulty of this generator, not
    real effect sizes.
    \item One Weibull shape per line, a shared coupling across arms, non-informative censoring,
    no outcome feedback across lines, and no treatment change within a line.
    \item Only $16$ covariates affect the outcome; the model sees many more real columns, which
    are non-causal here by construction.
    \item Arms are thin on late lines (\arm{ET alone} and \arm{ET+anti-CDK} on line 4): their
    counterfactuals are extrapolation, and unsupported arms are excluded from the metrics.
    \item Calendar time is an instrument and the hold-out enters late: an intentional shift
    that the reported errors include.
    \item The covariate standardiser is fitted on all samples, hold-out included (a property of
    the existing data module).
    \item Training settings were tuned on real data and chosen by hand here; the benchmark
    scores the model \emph{as trained}, not the best model it could be.
\end{enumerate}

% Needs only amsmath and booktabs, both already loaded there. Self-contained: no \todo,
% no external figures. Code references are to src/CausalSurv/semisynthetic/.

\providecommand{\drv}[1]{\texttt{#1}}
\providecommand{\arm}[1]{\textsc{#1}}
\providecommand{\Ind}{\mathbf{1}}

\section{Semi-Synthetic Data Generation}
\label{app:semisynthetic}

This appendix records every design choice of the semi-synthetic benchmark used for
counterfactual validation (Section~\ref{sec:counterfactual}): what was decided, and why.
The benchmark keeps the real covariates and treatment history of the cohort, and
\emph{simulates} the treatment of the line being scored, the potential survival time under
every arm, and the censoring. The true survival curve, restricted mean survival time (RMST)
and best arm of every patient under every arm are then known in closed form.

The numbers quoted below are for the reference configuration (assignment strength
$\gamma=1$, hidden-confounder strength $s=0$, heterogeneity $h=1$). All coefficients are
set by hand, not estimated: they fix the \emph{difficulty} of the benchmark, not its realism.
DynaSurv's own results on the sweep (Section~\ref{app:semisynthetic:sweep}) are reported
against the oracle and naive Kaplan--Meier references only: this rebuild does not yet include
the Cox/DeepSurv baselines used for the observed-vs-hidden-confounding argument in
Section~\ref{sec:counterfactual}.

% -----------------------------------------------------------------------------------------
\subsection{Cohort and arms}
\label{app:semisynthetic:cohort}

\paragraph{Cohort.} HR$^+$/HER2$^-$ patients whose first treatment line starts in 2018 or
later, lines 1--4, built exactly as the training pipeline builds its cohort (patient-level
entry-year filter, inner join with the static file), so that simulated rows line up one to
one with the rows the model trains on. This gives $6{,}191$ patients and $13{,}116$
(patient, line) records, $6{,}191 / 3{,}401 / 2{,}169 / 1{,}355$ on lines 1 to 4. The year
restriction is the identifiability window of the real pipeline: every arm was available
throughout. One exact duplicate (patient, line) record in the source file is removed, because
after expansion it would become a sample of its own with a false ``previous line''.
Administrative censoring is measured from the last observed line start, 2024-02-21.

\paragraph{Arms.} Four arms: \arm{ET alone}, \arm{ET+anti-CDK wo CT}, \arm{MonoCT std alone},
\arm{PolyCT alone}. \arm{CT+anti-angio} is not simulated (21, 7, 2 and 2 records per line in the
restricted cohort: no overlap). It and all other categories (\arm{CT+anti-HER2}, \arm{CT+IT},
\arm{CT+TT}, \arm{ET+TT}, \arm{No treatment}, \arm{Other}) remain in the data as
\emph{history} only and are excluded from the set of recommendable arms. Arm indices follow
the alphabetical order of the arm names throughout.

\paragraph{Real history, simulated present.} Covariates, past treatments and line structure
(previous-line duration, gap, line number) are real. Only the arm and outcome of the line
being scored are simulated, and the simulated outcome does \emph{not} feed back into later
lines. This keeps the covariate distribution real and makes each line a separate problem
given its history. The price is that the benchmark does not test how a model handles
treatment--outcome feedback across lines.

% -----------------------------------------------------------------------------------------
\subsection{Prefix expansion}
\label{app:semisynthetic:expansion}

A patient with $L$ lines becomes $L$ samples. Sample $j\in\{1,\dots,L\}$ contains the real
records of lines $1,\dots,j-1$ and record $j$ with the simulated arm and simulated outcome;
its identifier is $10\cdot\text{id}+j$, and all samples of a patient share the original id
(the unit of any train/validation grouping) and the entry year. The reference cohort yields
$13{,}116$ samples in $24{,}920$ rows.

The reason is that the sequence model steps through lines in order, and its hidden state at
line $j$ does not depend on the arm at line $j$. Scoring only the last line of each sample
therefore requires no change to the model or its losses; the alternative, simulating all
lines of a patient jointly, would require outcome feedback and a joint model. The data
module masks every row except the sample's last line; history rows carry placeholder
outcomes ($0$).

Columns describing the real current line and its outcome (drug flags other than the treatment
category, all outcome columns other than the two survival targets, line end date, death date)
are removed, so that no real value sits next to a simulated one. Static features do not reach
the encoder, so age and menopausal status are copied into dynamic columns; otherwise they
would act as unplanned hidden confounders.

% -----------------------------------------------------------------------------------------
\subsection{Drivers and causal roles}
\label{app:semisynthetic:drivers}

Each sample has a vector $z\in\mathbb{R}^{16}$ of drivers computed from its real record.
Continuous drivers are standardised over the cohort; the others are $0/1$ flags.
Table~\ref{tab:semisynthetic:drivers} gives their definitions and roles. The role of each
driver is enforced when the configuration is loaded: a coefficient on an instrument in the
outcome equation, or on a prognostic factor in the assignment equation, is rejected.

\begin{table}[ht]
\centering\small
\caption{Drivers of the data-generating process. ``Both'' means the driver enters assignment
and outcome (confounder).}
\label{tab:semisynthetic:drivers}
\begin{tabular}{@{}lll@{}}
\toprule
Driver & Definition & Role \\
\midrule
\drv{liver} & liver metastases $>0$ & both \\
\drv{visceral} & pulmonary or pleural metastases $>0$ & both \\
\drv{bone\_only} & bone metastases, none of liver, visceral, brain & both \\
\drv{mpps} & imputed ECOG/WHO performance status, 0--4 (standardised) & both \\
\drv{log\_prev\_line} & $\log(1+\text{previous line duration})$ (standardised) & both \\
\drv{prev\_cdk}, \drv{prev\_et}, \drv{prev\_ct} & class of the \emph{real} previous arm (one-hot; all zero at line 1) & both \\
\drv{cum\_new\_sites} & cumulative new metastatic sites (standardised) & both \\
\drv{age}, \drv{menopause} & age at selection (standardised), menopausal status & both \\
\drv{calendar} & months since 2018-01-01 (standardised) & assignment only \\
\drv{brain} & brain metastases $>0$ & outcome only \\
\drv{lobular} & invasive lobular histology & outcome only \\
\drv{brca} & \textit{BRCA1} or \textit{BRCA2} & outcome only \\
\drv{old\_site\_progression} & progression of an old site $>0$ & outcome only \\
\bottomrule
\end{tabular}
\end{table}

The model still receives all real covariate columns ($73$, plus the calendar covariate); only
these $16$ matter for the simulated data, and the remainder are nuisance columns it must
learn to ignore. A buffer-time covariate is deliberately not a driver.

\paragraph{Hidden confounder.} A latent variable
$u=\rho\,\tilde z_{\text{liver}}+\sqrt{1-\rho^2}\,\varepsilon$ with $\varepsilon\sim\mathcal{N}(0,1)$,
$\tilde z_{\text{liver}}$ the standardised liver indicator and $\rho=0.3$, enters both
assignment and outcome and is \emph{never} written to the model input. Correlating it with a
real covariate keeps it from being independent noise. Its outcome coefficient is fixed and
only its coupling to assignment is swept, so that strength $0$ is an unconfounded reference
carrying the same outcome heterogeneity: any error growth above it is confounding bias.

% -----------------------------------------------------------------------------------------
\subsection{Treatment assignment}
\label{app:semisynthetic:assignment}

For a sample on line $\ell$ with drivers $z$ and hidden $u$, the arm $A$ is drawn with
\begin{equation}
    \pi_a(z,u)=\mathbb{P}(A=a\mid z,u)
    =\frac{\exp\eta^{\text{A}}_a}{\sum_{b}\exp\eta^{\text{A}}_b},
    \qquad
    \eta^{\text{A}}_a=\alpha_{\ell a}+\gamma\,B_a^{\top}z+s\,c_a\,u,
    \label{eq:semisynthetic:assignment}
\end{equation}
where $B_a$ are hand-set coefficients (Table~\ref{tab:semisynthetic:coefficients}, shared
across lines), $c_a$ the hidden loadings, $\gamma$ the assignment strength, and $s$ the hidden
strength. Coefficients are chosen for clinically plausible signs (visceral involvement pushes
towards combination chemotherapy, a long previous line towards endocrine-based therapy, a CDK
inhibitor is rarely repeated) and moderate size.

\paragraph{Calibrated intercepts.} The intercepts $\alpha_{\ell a}$ are not set by hand. For
every line they are fitted, by iterative proportional fitting
\begin{equation}
    \alpha_{\ell a}\leftarrow\alpha_{\ell a}+\log\frac{q_{\ell a}}{\bar\pi_{\ell a}},
    \qquad
    \bar\pi_{\ell a}=\frac{1}{n_\ell}\sum_{i\in\ell}\pi_a(z_i,u_i),
\end{equation}
until $\max_a|\bar\pi_{\ell a}-q_{\ell a}|<10^{-6}$, where $q_{\ell\cdot}$ is the real
distribution of the four arms on line $\ell$. The marginal arm mix therefore equals the real
mix at every $\gamma$ and $s$: a sweep changes \emph{who} receives which arm, never \emph{how
many}. The average runs over every sample of the line (each is given a simulated arm) while
the target is the real mix among the four arms only.
The targets $q_\ell$ (order: ET, ET+CDK, MonoCT, PolyCT) are
$(.183,.628,.090,.100)$, $(.148,.412,.330,.109)$, $(.049,.107,.611,.233)$ and
$(.032,.031,.580,.357)$ for lines 1 to 4.

\paragraph{Sampling.} The arm is drawn by inverting the cumulative distribution with
\emph{one uniform per sample}, so that draws coincide wherever propensities coincide. The true
propensities are stored.

\paragraph{Range of the sweep.} The effective sample size of inverse-propensity weights of the
assigned arm, divided by $n$, is $0.51$, $0.40$, $0.22$, $0.01$ and $0.00$ at $\gamma=0$, $0.5$,
$1$, $2$ and $4$ respectively (no hidden confounder), and $0.07$ at $s=2$ even with $\gamma=0$. Beyond $\gamma\approx1.5$ or
$s\approx1$ the benchmark measures extrapolation rather than confounding, so the sweep is
restricted to $\gamma\le1.5$ and $s\le1$.

\begin{table}[ht]
\centering\small
\caption{Hand-set coefficients (blank $=0$). Left: assignment coefficients $B_a$ and hidden
loadings $c_a$. Right: outcome prognostic term $f$ and arm interactions $g_a$. Continuous
drivers are per standard deviation, flags are on \emph{vs} off; outcome coefficients are on the
log-hazard scale (positive $=$ worse). The row ``Hidden'' gives the loadings $c_a$ of $u$ on the
arms' assignment logits and, in the $f$ column, its outcome coefficient $\beta_u$.}
\label{tab:semisynthetic:coefficients}
\setlength{\tabcolsep}{4pt}
\begin{tabular}{@{}l rrrr c rrrrr@{}}
\toprule
 & \multicolumn{4}{c}{Assignment $B_a$} & & \multicolumn{5}{c}{Outcome} \\
\cmidrule(lr){2-5}\cmidrule(l){7-11}
Driver & ET & CDK & Mono & Poly & & $f$ & $g_{\text{ET}}$ & $g_{\text{CDK}}$ & $g_{\text{Mono}}$ & $g_{\text{Poly}}$ \\
\midrule
\drv{liver}          &      & $-0.3$ & $0.5$  & $0.9$  & & $0.45$ & $0.30$ & $0.25$ & $-0.15$ & $-0.35$ \\
\drv{visceral}       &      & $-0.2$ & $0.3$  & $0.5$  & & $0.25$ &        &        &         & $-0.25$ \\
\drv{bone\_only}     &      & $0.5$  & $-0.4$ & $-0.6$ & & $-0.30$ & $-0.15$ & $-0.30$ &      &         \\
\drv{mpps}           &      & $-0.3$ & $0.3$  & $-0.2$ & & $0.30$ &        &        & $-0.10$ & $0.15$  \\
\drv{log\_prev\_line} &     & $0.4$  & $-0.4$ & $-0.5$ & & $-0.25$ &       & $-0.20$ &        &         \\
\drv{prev\_cdk}      & $-0.3$ & $-1.0$ & $0.4$ & $0.3$  & & $0.20$ &        & $0.40$ &         &         \\
\drv{prev\_et}       & $-0.5$ &      &        &        & &        & $0.30$ &        &         &         \\
\drv{prev\_ct}       &      &        & $0.3$  &        & &        &        &        &         &         \\
\drv{cum\_new\_sites} &     & $-0.1$ & $0.2$  & $0.3$  & & $0.15$ &        &        &         &         \\
\drv{age}            & $0.4$ &       & $0.3$  & $-0.4$ & & $0.10$ &        &        & $-0.10$ & $0.15$  \\
\drv{menopause}      &      & $0.3$  &        & $-0.3$ & & $0.05$ &        &        &         &         \\
\drv{calendar}       &      & $0.3$  & $-0.1$ & $-0.2$ & &        &        &        &         &         \\
\drv{brain}          &      &        &        &        & & $0.40$ &        &        &         &         \\
\drv{lobular}        &      &        &        &        & & $0.10$ &        &        &         &         \\
\drv{brca}           &      &        &        &        & & $0.15$ &        &        &         &         \\
\drv{old\_site\_prog.} &    &        &        &        & & $0.15$ &        &        &         &         \\
\midrule
Hidden ($c_a$; $\beta_u$ under $f$) & 0 & $-0.5$ & $0.5$ & $1.0$ & & $0.4$ & & & & \\
Arm effect $\tau_a$  &      &        &        &        & & & $0$ & $-0.20$ & $0.10$ & $0$ \\
\bottomrule
\end{tabular}
\end{table}

% -----------------------------------------------------------------------------------------
\subsection{Potential outcomes}
\label{app:semisynthetic:outcome}

Time from the start of the line to death under arm $a$ follows a Weibull model with
line-specific shape $k_\ell$ and scale $\lambda_\ell$,
\begin{equation}
    S_a(t\mid z,u)=\exp\!\Big(-\big(t/\lambda_\ell\big)^{k_\ell}e^{\eta_a}\Big),
    \qquad
    \eta_a=f(z)+\tau_a+h\,g_a(z)+\beta_u\,u,
    \label{eq:semisynthetic:outcome}
\end{equation}
that is, $T_a\sim\text{Weibull}(k_\ell,\lambda_\ell e^{-\eta_a/k_\ell})$; $\eta_a>0$ is a higher
hazard. $f$ is the prognostic term, $\tau_a$ the arm effect, $g_a$ the interaction that makes
the best arm differ across patients (who benefits from what), $h$ the heterogeneity, and
$\beta_u=0.4$ the fixed outcome coefficient of the hidden confounder. The functions $f$ and
$g_a$ are linear in $z$ (Table~\ref{tab:semisynthetic:coefficients}) and \emph{centred within
each line}, which makes $(k_\ell,\lambda_\ell)$ the baseline of the average patient.

\paragraph{Fitting the baseline.} For each line, $(k_\ell,\lambda_\ell)$ minimise the squared
difference between the population survival implied by the simulation for the assigned arms,
$\bar S_\ell(t)=n_\ell^{-1}\sum_{i\in\ell}\exp\!\big(-(t/\lambda)^{k}e^{\eta_{i,A_i}}\big)$, and the
real per-line Kaplan--Meier curve, over the observed jump times at which at least $5\%$ of the
line is still at risk. The result is $k_\ell\approx1.27$--$1.33$ and
$\lambda_\ell\approx53/27/22/18$ months; the fitted medians are $48.5/22.6/16.3/12.7$ months
against $47.7/22.8/16.5/12.4$ in the real data.

\paragraph{Closed-form truth.} Survival is given by \eqref{eq:semisynthetic:outcome}. With
$\sigma_a=\lambda_\ell e^{-\eta_a/k_\ell}$ the RMST to horizon $\tau$ is
\begin{equation}
    \text{RMST}_a(\tau)=\int_0^{\tau}S_a(t)\,dt
    =\sigma_a\,\Gamma\!\big(1+\tfrac1{k_\ell}\big)\,P\!\big(\tfrac1{k_\ell},(\tau/\sigma_a)^{k_\ell}\big),
\end{equation}
with $P$ the regularised lower incomplete gamma function; the individual treatment effect of
arm $a$ against $b$ is $\text{RMST}_a-\text{RMST}_b$ and the best arm maximises RMST over the
recommendable arms.

\paragraph{Coupling of the arms.} A single uniform per sample is shared by all arms, so
$T_a=\sigma_a(-\log V)^{1/k_\ell}$ with the same $V\sim\mathcal{U}(0,1)$ for every $a$. The
arms then differ only through $\eta_a$, and the individual effect is not swamped by independent
noise. The marginals are unaffected; the joint distribution of potential outcomes is one
particular choice, and another coupling would give the same marginals with different
individual-level effects.

\paragraph{Effect sizes.} Effects are small on the RMST scale: the average true pairwise
differences at the reference configuration lie between $-0.44$ and $+0.42$ months, and the
\emph{headroom} of a policy (value of the oracle minus value of the best single arm per line)
is $0.00, 0.11, 0.47, 1.25, 1.71$ months for $h=0,0.5,1,2,3$. The arm effect of
\arm{ET+anti-CDK} was first $-0.35$, which made it the best arm for $63$--$74\%$ of patients
(by line) so that ``always CDK'' scored highly; at $-0.20$ its share on line 2 drops from
$63\%$ to $48\%$ and the best-arm metric discriminates.

% -----------------------------------------------------------------------------------------
\subsection{Censoring}
\label{app:semisynthetic:censoring}

The observed time and event indicator are $\min(T,C)$ and $\Ind\{T\le C\}$ with
\begin{equation}
    C=\min\big(C^{\text{adm}},\,C^{\text{drop}}\big),\qquad
    C^{\text{adm}}=\frac{t_{\text{cutoff}}-t^{\text{start}}}{30.44},\qquad
    C^{\text{drop}}\sim\text{Exp}(r_\ell),
\end{equation}
in months, with $T$ the latent time under the simulated arm. Administrative censoring is a
fact of the real cohort and is kept sample by sample. Dropout is added only where the cutoff
alone leaves the simulated event rate above the real one; $r_\ell$ is found by bisection on
$\log r$ (60 steps) so that the event rate equals the real rate of the restricted cohort,
using fixed unit-exponential draws so that the event rate is a deterministic decreasing
function of $r$. The targets are the event rates of the restricted cohort,
$0.442/0.547/0.578/0.596$ (not those of the full 2008-onwards cohort, $0.73$--$0.85$). With
the cutoff alone the rates would be $0.465/0.587/0.639/0.679$, so dropout is active on every
line, at $0.0024/0.0041/0.0073/0.0103$ per month.

Dropout draws are independent conditional on line and the fitted generator parameters.
Administrative follow-up depends on calendar date, which also affects assignment and is
associated with patient characteristics. Marginal independent censoring is therefore not
guaranteed, and training-cohort censoring weights need not transfer to the later test cohort.
Protocol v2 uses known survival probabilities and uncensored latent outcomes for primary
simulation scoring instead of making that transfer assumption.

% -----------------------------------------------------------------------------------------
\subsection{Randomness and artefacts}
\label{app:semisynthetic:seeds}

Four independent random streams (hidden confounder, assignment, outcome, censoring) are
spawned from one seed sequence keyed by (data seed, replicate). A sweep over $\gamma$, $s$ or
$h$ therefore reuses the same $u$, the same assignment uniforms, the same latent uniforms and
the same dropout draws: levels differ \emph{only} by the swept knob, and repeating a
configuration reproduces the data exactly.

Each replicate writes (i) the expanded dynamic and re-keyed static tables under the names of
the real data files, so that the directory can be used in place of the real data; (ii) a truth
table with, per sample, the assigned arm, hidden $u$, Weibull parameters, latent time under
every arm, both censoring times, observed time and event, and propensity and $\eta$ for every
arm; and (iii) a manifest with every parameter, the calibrated intercepts, the arm counts per
line, the fitted Weibull parameters and medians, the dropout rates, and the achieved and
target event rates. The hidden $u$ appears only in the truth table.

% -----------------------------------------------------------------------------------------
\subsection{Training and data-loading choices}
\label{app:semisynthetic:training}

\begin{itemize}
    \item \textbf{Mask.} Only the last line of each sample is unmasked, and the counts that
    decide which arms are recommendable are computed on those rows.
    \item \textbf{Grouped splits.} Every split other than the temporal one (random hold-out,
    cross-validation folds, early-stopping split) groups by original patient, because the
    samples of a patient share their history and would otherwise leak across the boundary. A
    temporal split already keeps them together (shared entry year).
    \item \textbf{Protocol-v2 validation.} Twenty percent of the pre-2021 original patients
    form a separate development validation set (fixed split seed 0); the remainder train
    the model, scaler and support models. The survival grid is fitted from training endpoint
    outcomes only. The temporal test set never supplies early stopping or hyperparameter
    selection. Historical results below predate this correction and require rerunning.
    \item \textbf{Temporal hold-out.} Patients entering in 2021 or later ($26\%$ of samples),
    a deliberate policy-era shift; calendar time is an instrument of the assignment, so the
    model receives it.
    \item \textbf{Support thresholds.} Arms need at least $30$ samples, $10$ events and $30$
    samples followed to the horizon at a line to be recommendable, against $200/30/100$ on
    real data, because the four arms thin out on later lines (on line 4, \arm{ET alone} has
    $25$ samples and is excluded there).
    \item \textbf{Optimisation.} Early stopping on the validation loss (patience $20$, at most
    $150$ epochs). The architecture and learning-rate settings are those tuned on the real
    data; they were not re-tuned for the synthetic data.
\end{itemize}

% -----------------------------------------------------------------------------------------
\subsection{Evaluation against the truth}
\label{app:semisynthetic:evaluation}

All quantities are computed on the temporal hold-out, at each sample's own line $\ell$ and its
horizon $\tau_\ell$ ($24, 18, 12, 12$ months), on a $0.1$-month grid, and only for arms that
are recommendable at that line. Write $\hat S_a$ for a predicted curve and $S_a$ for the truth.
\begin{itemize}
    \item \textbf{Curve error.} $\big(\tau_\ell^{-1}\!\int_0^{\tau_\ell}(\hat S_a-S_a)^2dt\big)^{1/2}$
    averaged in squares over samples, and the error and bias of $\widehat{\text{RMST}}_a-\text{RMST}_a$,
    reported separately for the assigned (factual) arm and for the other (counterfactual) arms.
    \item \textbf{Effect error.} For each arm pair $(a,b)$,
    $\text{PEHE}=\big(\mathbb{E}[(\hat\Delta_{ab}-\Delta_{ab})^2]\big)^{1/2}$ with
    $\Delta_{ab}=\text{RMST}_a-\text{RMST}_b$, the bias of the average effect
    $\mathbb{E}[\hat\Delta_{ab}]-\mathbb{E}[\Delta_{ab}]$, and the share of samples with the
    correct sign. Hidden confounding can affect both bias and PEHE.
    \item \textbf{Policy.} A policy picks the arm with the highest predicted RMST among
    supported arms (abstaining to the assigned arm where none is supported). Value is the
    true RMST of the chosen arm; regret is the oracle value minus it; the hit rate is the share
    of samples where the chosen arm is the true best arm. Reference policies: the assigned
    arm, the best single arm per line (an upper bound for constant policies), and uniform
    random.
    \item \textbf{Factual accuracy (protocol v2).} At the exact requested horizon, write
    $p=S_A(\tau_\ell)$ and $q=\hat S_A(\tau_\ell)$. The primary Brier score averages
    $p(1-q)^2+(1-p)q^2$. We also report Brier scores from uncensored latent outcomes and
    concordance of latent event times with horizon-specific predicted risk. These are
    simulation diagnostics, not estimators available on real censored data. The historical
    IPCW results below used training-cohort censoring weights and must be replaced.
    \item \textbf{References.} The \emph{oracle} (true curves; must score zero error and zero
    regret) and the \emph{naive Kaplan--Meier} (per line and arm, no covariates; the bar a
    covariate model must clear).
\end{itemize}

% -----------------------------------------------------------------------------------------
\subsection{Sweep design}
\label{app:semisynthetic:sweep}

One knob is varied at a time, the others staying at the reference values.
Table~\ref{tab:semisynthetic:sweep} lists the cells; each is replicated three times (data seed
and training seed equal to the replicate index), giving $30$ runs. The hidden-confounder axis
is the one that separates methods that adjust for observed covariates only from methods that
would be unbiased under hidden confounding; the assignment axis is expected to be flat for
outcome-regression estimators, since observed confounding adds variance rather than bias.

\begin{table}[ht]
\centering\small
\caption{Sweep cells. The reference cell is $\gamma=1$, $s=0$, $h=1$; $s=0$ and $h=1$ are
therefore not repeated.}
\label{tab:semisynthetic:sweep}
\begin{tabular}{@{}lll@{}}
\toprule
Knob & Levels & Meaning \\
\midrule
$\gamma$ & $0,\,0.25,\,0.5,\,1,\,1.5$ & assignment strength (observed confounding) \\
$s$ & $0.5,\,1$ (and $0$) & hidden confounder $\to$ assignment \\
$h$ & $0,\,0.5,\,2$ (and $1$) & effect heterogeneity \\
\bottomrule
\end{tabular}
\end{table}

% -----------------------------------------------------------------------------------------
\subsection{Sweep results}
\label{app:semisynthetic:sweep-results}

All $30$ runs of Table~\ref{tab:semisynthetic:sweep} completed and were scored on all four
checkpoint kinds (\texttt{val\_loss}, \texttt{bestCI}, \texttt{bestCALIB}, \texttt{final\_epoch});
Table~\ref{tab:semisynthetic:checkpoint} pools every cell and replicate by kind. \texttt{val\_loss}
--- the real-data recommender's default --- has the highest pooled regret and PEHE of the four;
\texttt{bestCI} is lowest on both and is used for Table~\ref{tab:semisynthetic:sweep-results}.
This was a post-hoc choice using the reported hold-out's oracle metrics, while that hold-out
also supplied early stopping. The following tables are therefore exploratory historical
results, not independent validation. Protocol-v2 aggregation accepts a fixed checkpoint rule
(default \texttt{val\_loss}) and never selects it from the reported oracle regret.

\paragraph{Comparator interpretation.} The omitted Kaplan--Meier effect-error comparison
changes the conclusion: at these historical \texttt{bestCI} checkpoints DynaSurv has higher
mean pair PEHE in nine of ten cells. The reference values are $1.761$ versus $1.705$ months;
at hidden strength 1 they are $2.699$ versus $2.171$; at heterogeneity 0, $2.082$ versus
$1.192$. Only heterogeneity 2 favors DynaSurv ($2.643$ versus $2.773$). Reference policy
regret favors DynaSurv ($0.336$ versus $0.516$ months). The corrected aggregator reports
curve, effect, policy and factual metrics for every comparator on the same support. Ranking
benefit must not be interpreted as accurate individual treatment-effect estimation.

\begin{table}[ht]
\centering\small
\caption{Checkpoint-kind comparison, pooled over all 30 sweep runs. PEHE and $|$bias$|$
(months) are $n$-weighted means over arm pairs and cells; regret (months) and hit rate are
$n$-weighted means of the \texttt{line\_only} policy; factual $C$ is the IPCW concordance index.}
\label{tab:semisynthetic:checkpoint}
\begin{tabular}{@{}lrrrrr@{}}
\toprule
Checkpoint & PEHE & $|$bias$|$ & Regret & Hit rate & Factual $C$ \\
\midrule
\texttt{val\_loss} (real-data default) & 2.37 & 1.51 & 0.42 & 0.62 & 0.666 \\
\texttt{bestCI} & 1.95 & 1.14 & 0.36 & 0.64 & 0.675 \\
\texttt{bestCALIB} & 2.43 & 1.35 & 0.37 & 0.58 & 0.665 \\
\texttt{final\_epoch} & 2.23 & 1.16 & 0.40 & 0.52 & 0.663 \\
\bottomrule
\end{tabular}
\end{table}

\begin{table}[ht]
\centering\small
\caption{Sweep results at the \texttt{bestCI} checkpoint, mean $\pm$ std over 3 replicates.
PEHE and $|$bias$|$ (months) pool over lines and arm pairs; regret (months, oracle $=0$) and
hit rate use the \texttt{line\_only} policy; $h_r$ is the random policy's regret, an upper
bound on the task's headroom at that cell.}
\label{tab:semisynthetic:sweep-results}
\begin{tabular}{@{}llrrrrr@{}}
\toprule
Axis & Level & PEHE & $|$bias$|$ & Regret & Hit rate & $h_r$ \\
\midrule
$\gamma$ & 0        & $1.27\pm0.15$ & $0.34\pm0.15$ & $0.41\pm0.03$ & $0.57\pm0.05$ & 0.77 \\
$\gamma$ & 0.25      & $1.45\pm0.28$ & $0.70\pm0.40$ & $0.45\pm0.02$ & $0.56\pm0.02$ & 0.78 \\
$\gamma$ & 0.5       & $1.61\pm0.19$ & $0.85\pm0.23$ & $0.39\pm0.04$ & $0.59\pm0.04$ & 0.79 \\
$\gamma$ & 1 (ref.)  & $1.76\pm0.42$ & $0.87\pm0.33$ & $0.34\pm0.04$ & $0.64\pm0.04$ & 0.80 \\
$\gamma$ & 1.5       & $2.33\pm1.10$ & $1.49\pm1.03$ & $0.40\pm0.05$ & $0.62\pm0.02$ & 0.81 \\
$s$ & 0.5            & $2.25\pm0.32$ & $1.49\pm0.24$ & $0.44\pm0.04$ & $0.59\pm0.02$ & 0.80 \\
$s$ & 1               & $2.70\pm0.15$ & $2.14\pm0.16$ & $0.48\pm0.01$ & $0.57\pm0.01$ & 0.80 \\
$h$ & 0               & $2.08\pm0.75$ & $1.28\pm0.34$ & $0.03\pm0.03$ & $0.95\pm0.03$ & 0.46 \\
$h$ & 0.5             & $1.46\pm0.29$ & $0.94\pm0.45$ & $0.15\pm0.08$ & $0.67\pm0.12$ & 0.50 \\
$h$ & 2               & $2.64\pm0.53$ & $1.35\pm0.49$ & $0.52\pm0.07$ & $0.62\pm0.01$ & 1.46 \\
\bottomrule
\end{tabular}
\end{table}

\begin{table}[ht]
\centering\small
\caption{Overlap (effective sample size / $n$ per line, mean over 3 replicates), from the
synced manifests. Heterogeneity does not touch assignment, so it repeats the $\gamma=1$ row
exactly at every level.}
\label{tab:semisynthetic:overlap}
\begin{tabular}{@{}llrrrr@{}}
\toprule
Axis & Level & Line 1 & Line 2 & Line 3 & Line 4 \\
\midrule
$\gamma$ & 0        & 0.566 & 0.750 & 0.440 & 0.239 \\
$\gamma$ & 0.25      & 0.539 & 0.709 & 0.422 & 0.217 \\
$\gamma$ & 0.5       & 0.462 & 0.586 & 0.364 & 0.184 \\
$\gamma$ & 1 (ref.)  & 0.251 & 0.282 & 0.198 & 0.086 \\
$\gamma$ & 1.5       & 0.087 & 0.154 & 0.085 & 0.097 \\
$s$ & 0.5            & 0.197 & 0.275 & 0.140 & 0.061 \\
$s$ & 1               & 0.156 & 0.199 & 0.109 & 0.098 \\
$h$ & any level        & 0.251 & 0.282 & 0.198 & 0.086 \\
\bottomrule
\end{tabular}
\end{table}

\paragraph{Reading the three axes.} $\gamma$ (assignment on observed covariates) roughly
doubles PEHE and the average-effect bias from $\gamma=0$ to $\gamma=1.5$, but the model's own
regret is flat within noise (0.34--0.45) while the random-policy headroom $h_r$ grows only
slightly (0.77 to 0.81): the ranking of arms survives better than the size of the effect does.
This is consistent with overlap collapse: mean effective sample size / $n$ on line 4 falls from
$0.24$ at $\gamma=0$ to $0.10$ at $\gamma=1.5$ (Table~\ref{tab:semisynthetic:overlap}), so
precision, not direction, is what degrades. Three replicates cannot separate true bias from
finite-sample noise on this axis; the sweep gives a variance estimate, not a significance test.
$s$ (hidden confounder into assignment) is the sharper effect: bias more than doubles from the
$s=0$ reference ($0.87$) to $s=1$ ($2.14$) and regret rises by $40\%$ even though its overlap
loss is comparable to $\gamma$'s (line-4 ESS/$n$ falls to $0.06$--$0.10$) -- the pattern expected
of a confounder the model cannot see, not merely thinner support. $h$ (heterogeneity) mainly
changes the size of the task rather than the model's absolute accuracy: it does not change
assignment at all (identical overlap to the $\gamma=1$ reference row, Table~\ref{tab:semisynthetic:overlap}),
so any shift in the metrics below comes from the outcome model alone. At $h=0$ every line has
a single best arm, so both the model and a constant policy score near zero regret ($h_r=0.46$
is administrative-censoring headroom, not heterogeneity); as $h$ grows to $2$, $h_r$ triples to
$1.46$ and the model's regret grows with it (to $0.52$) while staying well under the random
policy's ceiling and the hit rate falls from $0.95$ to $0.62$ -- picking the right arm gets
harder as there are more genuinely different arms to pick between.

\paragraph{What this does not show.} The 30-run sweep uses one training configuration
(Section~\ref{app:semisynthetic:training}) and the model config predating this benchmark; it is
not re-tuned per cell. It has no Cox or DeepSurv baseline, so it cannot yet reproduce the
observed-vs-hidden-confounding bias comparison of Section~\ref{sec:counterfactual} with this
generator. Achieved censoring event rates matched their per-line targets to within $0.0007$ in
every one of the 30 synced manifests, confirming the calibration of
Section~\ref{app:semisynthetic:censoring} generalizes across the sweep.

% -----------------------------------------------------------------------------------------
\subsection{Verification of the generator}
\label{app:semisynthetic:verification}

The following were checked on the reference replicate.
\begin{itemize}
    \item Mean propensity per line equals the target to $10^{-6}$ for
    $\gamma\in\{0,0.5,1,2,4\}$ and $s\in\{0,2\}$; sampled arm counts pass a chi-square test;
    at $\gamma=s=0$ the propensity is constant within a line.
    \item A multinomial regression of the sampled arm on the drivers, pooled over lines with
    line indicators, recovers the assignment coefficients (correlation $0.976$, all signs
    correct). On line 1 alone it does not ($0.58$), because the previous-line drivers are
    constant there.
    \item Simulated Kaplan--Meier curves match the real ones (within about $0.04$ at 6, 12,
    24 and 36 months, medians within $2\%$); observed curves after censoring within about
    $0.02$; simulated event rates equal the real rates.
    \item Closed-form RMST agrees with Monte Carlo (bias $-0.002$ months, standardised
    differences of standard deviation $1.01$); at $h=0$, $\tau_a=0$ all arms have equal $\eta$.
    \item Expansion: the unmasked rows per line equal the real patients per line; history rows
    equal the real records; no hidden variable and no propensity in the model-facing files.
    \item Reproducibility: an identical rerun is bit-identical; changing $\gamma$ keeps $u$
    and changes $13\%$ of the arms; a new replicate changes $u$.
    \item No patient appears on both sides of a temporal, random or cross-validation split.
\end{itemize}

% -----------------------------------------------------------------------------------------
\subsection{Limitations}
\label{app:semisynthetic:limitations}

\begin{enumerate}
    \item Coefficients are hand-set; results describe the difficulty of this generator, not
    real effect sizes.
    \item One Weibull shape per line, a shared coupling across arms, non-informative censoring,
    no outcome feedback across lines, and no treatment change within a line.
    \item Only $16$ covariates affect the outcome; the model sees many more real columns, which
    are non-causal here by construction.
    \item Arms are thin on late lines (\arm{ET alone} and \arm{ET+anti-CDK} on line 4): their
    counterfactuals are extrapolation, and unsupported arms are excluded from the metrics.
    \item Calendar time is an instrument and the hold-out enters late: an intentional shift
    that the reported errors include.
    \item The covariate standardiser is fitted on all samples, hold-out included (a property of
    the existing data module).
    \item Training settings were tuned on real data and chosen by hand here; the benchmark
    scores the model \emph{as trained}, not the best model it could be.
\end{enumerate}

% Needs only amsmath and booktabs, both already loaded there. Self-contained: no \todo,
% no external figures. Code references are to src/CausalSurv/semisynthetic/.

\providecommand{\drv}[1]{\texttt{#1}}
\providecommand{\arm}[1]{\textsc{#1}}
\providecommand{\Ind}{\mathbf{1}}

\section{Semi-Synthetic Data Generation}
\label{app:semisynthetic}

This appendix records every design choice of the semi-synthetic benchmark used for
counterfactual validation (Section~\ref{sec:counterfactual}): what was decided, and why.
The benchmark keeps the real covariates and treatment history of the cohort, and
\emph{simulates} the treatment of the line being scored, the potential survival time under
every arm, and the censoring. The true survival curve, restricted mean survival time (RMST)
and best arm of every patient under every arm are then known in closed form.

The numbers quoted below are for the reference configuration (assignment strength
$\gamma=1$, hidden-confounder strength $s=0$, heterogeneity $h=1$). All coefficients are
set by hand, not estimated: they fix the \emph{difficulty} of the benchmark, not its realism.
DynaSurv's own results on the sweep (Section~\ref{app:semisynthetic:sweep}) are reported
against the oracle and naive Kaplan--Meier references only: this rebuild does not yet include
the Cox/DeepSurv baselines used for the observed-vs-hidden-confounding argument in
Section~\ref{sec:counterfactual}.

% -----------------------------------------------------------------------------------------
\subsection{Cohort and arms}
\label{app:semisynthetic:cohort}

\paragraph{Cohort.} HR$^+$/HER2$^-$ patients whose first treatment line starts in 2018 or
later, lines 1--4, built exactly as the training pipeline builds its cohort (patient-level
entry-year filter, inner join with the static file), so that simulated rows line up one to
one with the rows the model trains on. This gives $6{,}191$ patients and $13{,}116$
(patient, line) records, $6{,}191 / 3{,}401 / 2{,}169 / 1{,}355$ on lines 1 to 4. The year
restriction is the identifiability window of the real pipeline: every arm was available
throughout. One exact duplicate (patient, line) record in the source file is removed, because
after expansion it would become a sample of its own with a false ``previous line''.
Administrative censoring is measured from the last observed line start, 2024-02-21.

\paragraph{Arms.} Four arms: \arm{ET alone}, \arm{ET+anti-CDK wo CT}, \arm{MonoCT std alone},
\arm{PolyCT alone}. \arm{CT+anti-angio} is not simulated (21, 7, 2 and 2 records per line in the
restricted cohort: no overlap). It and all other categories (\arm{CT+anti-HER2}, \arm{CT+IT},
\arm{CT+TT}, \arm{ET+TT}, \arm{No treatment}, \arm{Other}) remain in the data as
\emph{history} only and are excluded from the set of recommendable arms. Arm indices follow
the alphabetical order of the arm names throughout.

\paragraph{Real history, simulated present.} Covariates, past treatments and line structure
(previous-line duration, gap, line number) are real. Only the arm and outcome of the line
being scored are simulated, and the simulated outcome does \emph{not} feed back into later
lines. This keeps the covariate distribution real and makes each line a separate problem
given its history. The price is that the benchmark does not test how a model handles
treatment--outcome feedback across lines.

% -----------------------------------------------------------------------------------------
\subsection{Prefix expansion}
\label{app:semisynthetic:expansion}

A patient with $L$ lines becomes $L$ samples. Sample $j\in\{1,\dots,L\}$ contains the real
records of lines $1,\dots,j-1$ and record $j$ with the simulated arm and simulated outcome;
its identifier is $10\cdot\text{id}+j$, and all samples of a patient share the original id
(the unit of any train/validation grouping) and the entry year. The reference cohort yields
$13{,}116$ samples in $24{,}920$ rows.

The reason is that the sequence model steps through lines in order, and its hidden state at
line $j$ does not depend on the arm at line $j$. Scoring only the last line of each sample
therefore requires no change to the model or its losses; the alternative, simulating all
lines of a patient jointly, would require outcome feedback and a joint model. The data
module masks every row except the sample's last line; history rows carry placeholder
outcomes ($0$).

Columns describing the real current line and its outcome (drug flags other than the treatment
category, all outcome columns other than the two survival targets, line end date, death date)
are removed, so that no real value sits next to a simulated one. Static features do not reach
the encoder, so age and menopausal status are copied into dynamic columns; otherwise they
would act as unplanned hidden confounders.

% -----------------------------------------------------------------------------------------
\subsection{Drivers and causal roles}
\label{app:semisynthetic:drivers}

Each sample has a vector $z\in\mathbb{R}^{16}$ of drivers computed from its real record.
Continuous drivers are standardised over the cohort; the others are $0/1$ flags.
Table~\ref{tab:semisynthetic:drivers} gives their definitions and roles. The role of each
driver is enforced when the configuration is loaded: a coefficient on an instrument in the
outcome equation, or on a prognostic factor in the assignment equation, is rejected.

\begin{table}[ht]
\centering\small
\caption{Drivers of the data-generating process. ``Both'' means the driver enters assignment
and outcome (confounder).}
\label{tab:semisynthetic:drivers}
\begin{tabular}{@{}lll@{}}
\toprule
Driver & Definition & Role \\
\midrule
\drv{liver} & liver metastases $>0$ & both \\
\drv{visceral} & pulmonary or pleural metastases $>0$ & both \\
\drv{bone\_only} & bone metastases, none of liver, visceral, brain & both \\
\drv{mpps} & imputed ECOG/WHO performance status, 0--4 (standardised) & both \\
\drv{log\_prev\_line} & $\log(1+\text{previous line duration})$ (standardised) & both \\
\drv{prev\_cdk}, \drv{prev\_et}, \drv{prev\_ct} & class of the \emph{real} previous arm (one-hot; all zero at line 1) & both \\
\drv{cum\_new\_sites} & cumulative new metastatic sites (standardised) & both \\
\drv{age}, \drv{menopause} & age at selection (standardised), menopausal status & both \\
\drv{calendar} & months since 2018-01-01 (standardised) & assignment only \\
\drv{brain} & brain metastases $>0$ & outcome only \\
\drv{lobular} & invasive lobular histology & outcome only \\
\drv{brca} & \textit{BRCA1} or \textit{BRCA2} & outcome only \\
\drv{old\_site\_progression} & progression of an old site $>0$ & outcome only \\
\bottomrule
\end{tabular}
\end{table}

The model still receives all real covariate columns ($73$, plus the calendar covariate); only
these $16$ matter for the simulated data, and the remainder are nuisance columns it must
learn to ignore. A buffer-time covariate is deliberately not a driver.

\paragraph{Hidden confounder.} A latent variable
$u=\rho\,\tilde z_{\text{liver}}+\sqrt{1-\rho^2}\,\varepsilon$ with $\varepsilon\sim\mathcal{N}(0,1)$,
$\tilde z_{\text{liver}}$ the standardised liver indicator and $\rho=0.3$, enters both
assignment and outcome and is \emph{never} written to the model input. Correlating it with a
real covariate keeps it from being independent noise. Its outcome coefficient is fixed and
only its coupling to assignment is swept, so that strength $0$ is an unconfounded reference
carrying the same outcome heterogeneity: any error growth above it is confounding bias.

% -----------------------------------------------------------------------------------------
\subsection{Treatment assignment}
\label{app:semisynthetic:assignment}

For a sample on line $\ell$ with drivers $z$ and hidden $u$, the arm $A$ is drawn with
\begin{equation}
    \pi_a(z,u)=\mathbb{P}(A=a\mid z,u)
    =\frac{\exp\eta^{\text{A}}_a}{\sum_{b}\exp\eta^{\text{A}}_b},
    \qquad
    \eta^{\text{A}}_a=\alpha_{\ell a}+\gamma\,B_a^{\top}z+s\,c_a\,u,
    \label{eq:semisynthetic:assignment}
\end{equation}
where $B_a$ are hand-set coefficients (Table~\ref{tab:semisynthetic:coefficients}, shared
across lines), $c_a$ the hidden loadings, $\gamma$ the assignment strength, and $s$ the hidden
strength. Coefficients are chosen for clinically plausible signs (visceral involvement pushes
towards combination chemotherapy, a long previous line towards endocrine-based therapy, a CDK
inhibitor is rarely repeated) and moderate size.

\paragraph{Calibrated intercepts.} The intercepts $\alpha_{\ell a}$ are not set by hand. For
every line they are fitted, by iterative proportional fitting
\begin{equation}
    \alpha_{\ell a}\leftarrow\alpha_{\ell a}+\log\frac{q_{\ell a}}{\bar\pi_{\ell a}},
    \qquad
    \bar\pi_{\ell a}=\frac{1}{n_\ell}\sum_{i\in\ell}\pi_a(z_i,u_i),
\end{equation}
until $\max_a|\bar\pi_{\ell a}-q_{\ell a}|<10^{-6}$, where $q_{\ell\cdot}$ is the real
distribution of the four arms on line $\ell$. The marginal arm mix therefore equals the real
mix at every $\gamma$ and $s$: a sweep changes \emph{who} receives which arm, never \emph{how
many}. The average runs over every sample of the line (each is given a simulated arm) while
the target is the real mix among the four arms only.
The targets $q_\ell$ (order: ET, ET+CDK, MonoCT, PolyCT) are
$(.183,.628,.090,.100)$, $(.148,.412,.330,.109)$, $(.049,.107,.611,.233)$ and
$(.032,.031,.580,.357)$ for lines 1 to 4.

\paragraph{Sampling.} The arm is drawn by inverting the cumulative distribution with
\emph{one uniform per sample}, so that draws coincide wherever propensities coincide. The true
propensities are stored.

\paragraph{Range of the sweep.} The effective sample size of inverse-propensity weights of the
assigned arm, divided by $n$, is $0.51$, $0.40$, $0.22$, $0.01$ and $0.00$ at $\gamma=0$, $0.5$,
$1$, $2$ and $4$ respectively (no hidden confounder), and $0.07$ at $s=2$ even with $\gamma=0$. Beyond $\gamma\approx1.5$ or
$s\approx1$ the benchmark measures extrapolation rather than confounding, so the sweep is
restricted to $\gamma\le1.5$ and $s\le1$.

\begin{table}[ht]
\centering\small
\caption{Hand-set coefficients (blank $=0$). Left: assignment coefficients $B_a$ and hidden
loadings $c_a$. Right: outcome prognostic term $f$ and arm interactions $g_a$. Continuous
drivers are per standard deviation, flags are on \emph{vs} off; outcome coefficients are on the
log-hazard scale (positive $=$ worse). The row ``Hidden'' gives the loadings $c_a$ of $u$ on the
arms' assignment logits and, in the $f$ column, its outcome coefficient $\beta_u$.}
\label{tab:semisynthetic:coefficients}
\setlength{\tabcolsep}{4pt}
\begin{tabular}{@{}l rrrr c rrrrr@{}}
\toprule
 & \multicolumn{4}{c}{Assignment $B_a$} & & \multicolumn{5}{c}{Outcome} \\
\cmidrule(lr){2-5}\cmidrule(l){7-11}
Driver & ET & CDK & Mono & Poly & & $f$ & $g_{\text{ET}}$ & $g_{\text{CDK}}$ & $g_{\text{Mono}}$ & $g_{\text{Poly}}$ \\
\midrule
\drv{liver}          &      & $-0.3$ & $0.5$  & $0.9$  & & $0.45$ & $0.30$ & $0.25$ & $-0.15$ & $-0.35$ \\
\drv{visceral}       &      & $-0.2$ & $0.3$  & $0.5$  & & $0.25$ &        &        &         & $-0.25$ \\
\drv{bone\_only}     &      & $0.5$  & $-0.4$ & $-0.6$ & & $-0.30$ & $-0.15$ & $-0.30$ &      &         \\
\drv{mpps}           &      & $-0.3$ & $0.3$  & $-0.2$ & & $0.30$ &        &        & $-0.10$ & $0.15$  \\
\drv{log\_prev\_line} &     & $0.4$  & $-0.4$ & $-0.5$ & & $-0.25$ &       & $-0.20$ &        &         \\
\drv{prev\_cdk}      & $-0.3$ & $-1.0$ & $0.4$ & $0.3$  & & $0.20$ &        & $0.40$ &         &         \\
\drv{prev\_et}       & $-0.5$ &      &        &        & &        & $0.30$ &        &         &         \\
\drv{prev\_ct}       &      &        & $0.3$  &        & &        &        &        &         &         \\
\drv{cum\_new\_sites} &     & $-0.1$ & $0.2$  & $0.3$  & & $0.15$ &        &        &         &         \\
\drv{age}            & $0.4$ &       & $0.3$  & $-0.4$ & & $0.10$ &        &        & $-0.10$ & $0.15$  \\
\drv{menopause}      &      & $0.3$  &        & $-0.3$ & & $0.05$ &        &        &         &         \\
\drv{calendar}       &      & $0.3$  & $-0.1$ & $-0.2$ & &        &        &        &         &         \\
\drv{brain}          &      &        &        &        & & $0.40$ &        &        &         &         \\
\drv{lobular}        &      &        &        &        & & $0.10$ &        &        &         &         \\
\drv{brca}           &      &        &        &        & & $0.15$ &        &        &         &         \\
\drv{old\_site\_prog.} &    &        &        &        & & $0.15$ &        &        &         &         \\
\midrule
Hidden ($c_a$; $\beta_u$ under $f$) & 0 & $-0.5$ & $0.5$ & $1.0$ & & $0.4$ & & & & \\
Arm effect $\tau_a$  &      &        &        &        & & & $0$ & $-0.20$ & $0.10$ & $0$ \\
\bottomrule
\end{tabular}
\end{table}

% -----------------------------------------------------------------------------------------
\subsection{Potential outcomes}
\label{app:semisynthetic:outcome}

Time from the start of the line to death under arm $a$ follows a Weibull model with
line-specific shape $k_\ell$ and scale $\lambda_\ell$,
\begin{equation}
    S_a(t\mid z,u)=\exp\!\Big(-\big(t/\lambda_\ell\big)^{k_\ell}e^{\eta_a}\Big),
    \qquad
    \eta_a=f(z)+\tau_a+h\,g_a(z)+\beta_u\,u,
    \label{eq:semisynthetic:outcome}
\end{equation}
that is, $T_a\sim\text{Weibull}(k_\ell,\lambda_\ell e^{-\eta_a/k_\ell})$; $\eta_a>0$ is a higher
hazard. $f$ is the prognostic term, $\tau_a$ the arm effect, $g_a$ the interaction that makes
the best arm differ across patients (who benefits from what), $h$ the heterogeneity, and
$\beta_u=0.4$ the fixed outcome coefficient of the hidden confounder. The functions $f$ and
$g_a$ are linear in $z$ (Table~\ref{tab:semisynthetic:coefficients}) and \emph{centred within
each line}, which makes $(k_\ell,\lambda_\ell)$ the baseline of the average patient.

\paragraph{Fitting the baseline.} For each line, $(k_\ell,\lambda_\ell)$ minimise the squared
difference between the population survival implied by the simulation for the assigned arms,
$\bar S_\ell(t)=n_\ell^{-1}\sum_{i\in\ell}\exp\!\big(-(t/\lambda)^{k}e^{\eta_{i,A_i}}\big)$, and the
real per-line Kaplan--Meier curve, over the observed jump times at which at least $5\%$ of the
line is still at risk. The result is $k_\ell\approx1.27$--$1.33$ and
$\lambda_\ell\approx53/27/22/18$ months; the fitted medians are $48.5/22.6/16.3/12.7$ months
against $47.7/22.8/16.5/12.4$ in the real data.

\paragraph{Closed-form truth.} Survival is given by \eqref{eq:semisynthetic:outcome}. With
$\sigma_a=\lambda_\ell e^{-\eta_a/k_\ell}$ the RMST to horizon $\tau$ is
\begin{equation}
    \text{RMST}_a(\tau)=\int_0^{\tau}S_a(t)\,dt
    =\sigma_a\,\Gamma\!\big(1+\tfrac1{k_\ell}\big)\,P\!\big(\tfrac1{k_\ell},(\tau/\sigma_a)^{k_\ell}\big),
\end{equation}
with $P$ the regularised lower incomplete gamma function; the individual treatment effect of
arm $a$ against $b$ is $\text{RMST}_a-\text{RMST}_b$ and the best arm maximises RMST over the
recommendable arms.

\paragraph{Coupling of the arms.} A single uniform per sample is shared by all arms, so
$T_a=\sigma_a(-\log V)^{1/k_\ell}$ with the same $V\sim\mathcal{U}(0,1)$ for every $a$. The
arms then differ only through $\eta_a$, and the individual effect is not swamped by independent
noise. The marginals are unaffected; the joint distribution of potential outcomes is one
particular choice, and another coupling would give the same marginals with different
individual-level effects.

\paragraph{Effect sizes.} Effects are small on the RMST scale: the average true pairwise
differences at the reference configuration lie between $-0.44$ and $+0.42$ months, and the
\emph{headroom} of a policy (value of the oracle minus value of the best single arm per line)
is $0.00, 0.11, 0.47, 1.25, 1.71$ months for $h=0,0.5,1,2,3$. The arm effect of
\arm{ET+anti-CDK} was first $-0.35$, which made it the best arm for $63$--$74\%$ of patients
(by line) so that ``always CDK'' scored highly; at $-0.20$ its share on line 2 drops from
$63\%$ to $48\%$ and the best-arm metric discriminates.

% -----------------------------------------------------------------------------------------
\subsection{Censoring}
\label{app:semisynthetic:censoring}

The observed time and event indicator are $\min(T,C)$ and $\Ind\{T\le C\}$ with
\begin{equation}
    C=\min\big(C^{\text{adm}},\,C^{\text{drop}}\big),\qquad
    C^{\text{adm}}=\frac{t_{\text{cutoff}}-t^{\text{start}}}{30.44},\qquad
    C^{\text{drop}}\sim\text{Exp}(r_\ell),
\end{equation}
in months, with $T$ the latent time under the simulated arm. Administrative censoring is a
fact of the real cohort and is kept sample by sample. Dropout is added only where the cutoff
alone leaves the simulated event rate above the real one; $r_\ell$ is found by bisection on
$\log r$ (60 steps) so that the event rate equals the real rate of the restricted cohort,
using fixed unit-exponential draws so that the event rate is a deterministic decreasing
function of $r$. The targets are the event rates of the restricted cohort,
$0.442/0.547/0.578/0.596$ (not those of the full 2008-onwards cohort, $0.73$--$0.85$). With
the cutoff alone the rates would be $0.465/0.587/0.639/0.679$, so dropout is active on every
line, at $0.0024/0.0041/0.0073/0.0103$ per month.

Dropout draws are independent conditional on line and the fitted generator parameters.
Administrative follow-up depends on calendar date, which also affects assignment and is
associated with patient characteristics. Marginal independent censoring is therefore not
guaranteed, and training-cohort censoring weights need not transfer to the later test cohort.
Protocol v2 uses known survival probabilities and uncensored latent outcomes for primary
simulation scoring instead of making that transfer assumption.

% -----------------------------------------------------------------------------------------
\subsection{Randomness and artefacts}
\label{app:semisynthetic:seeds}

Four independent random streams (hidden confounder, assignment, outcome, censoring) are
spawned from one seed sequence keyed by (data seed, replicate). A sweep over $\gamma$, $s$ or
$h$ therefore reuses the same $u$, the same assignment uniforms, the same latent uniforms and
the same dropout draws: levels differ \emph{only} by the swept knob, and repeating a
configuration reproduces the data exactly.

Each replicate writes (i) the expanded dynamic and re-keyed static tables under the names of
the real data files, so that the directory can be used in place of the real data; (ii) a truth
table with, per sample, the assigned arm, hidden $u$, Weibull parameters, latent time under
every arm, both censoring times, observed time and event, and propensity and $\eta$ for every
arm; and (iii) a manifest with every parameter, the calibrated intercepts, the arm counts per
line, the fitted Weibull parameters and medians, the dropout rates, and the achieved and
target event rates. The hidden $u$ appears only in the truth table.

% -----------------------------------------------------------------------------------------
\subsection{Training and data-loading choices}
\label{app:semisynthetic:training}

\begin{itemize}
    \item \textbf{Mask.} Only the last line of each sample is unmasked, and the counts that
    decide which arms are recommendable are computed on those rows.
    \item \textbf{Grouped splits.} Every split other than the temporal one (random hold-out,
    cross-validation folds, early-stopping split) groups by original patient, because the
    samples of a patient share their history and would otherwise leak across the boundary. A
    temporal split already keeps them together (shared entry year).
    \item \textbf{Protocol-v2 validation.} Twenty percent of the pre-2021 original patients
    form a separate development validation set (fixed split seed 0); the remainder train
    the model, scaler and support models. The survival grid is fitted from training endpoint
    outcomes only. The temporal test set never supplies early stopping or hyperparameter
    selection. Historical results below predate this correction and require rerunning.
    \item \textbf{Temporal hold-out.} Patients entering in 2021 or later ($26\%$ of samples),
    a deliberate policy-era shift; calendar time is an instrument of the assignment, so the
    model receives it.
    \item \textbf{Support thresholds.} Arms need at least $30$ samples, $10$ events and $30$
    samples followed to the horizon at a line to be recommendable, against $200/30/100$ on
    real data, because the four arms thin out on later lines (on line 4, \arm{ET alone} has
    $25$ samples and is excluded there).
    \item \textbf{Optimisation.} Early stopping on the validation loss (patience $20$, at most
    $150$ epochs). The architecture and learning-rate settings are those tuned on the real
    data; they were not re-tuned for the synthetic data.
\end{itemize}

% -----------------------------------------------------------------------------------------
\subsection{Evaluation against the truth}
\label{app:semisynthetic:evaluation}

All quantities are computed on the temporal hold-out, at each sample's own line $\ell$ and its
horizon $\tau_\ell$ ($24, 18, 12, 12$ months), on a $0.1$-month grid, and only for arms that
are recommendable at that line. Write $\hat S_a$ for a predicted curve and $S_a$ for the truth.
\begin{itemize}
    \item \textbf{Curve error.} $\big(\tau_\ell^{-1}\!\int_0^{\tau_\ell}(\hat S_a-S_a)^2dt\big)^{1/2}$
    averaged in squares over samples, and the error and bias of $\widehat{\text{RMST}}_a-\text{RMST}_a$,
    reported separately for the assigned (factual) arm and for the other (counterfactual) arms.
    \item \textbf{Effect error.} For each arm pair $(a,b)$,
    $\text{PEHE}=\big(\mathbb{E}[(\hat\Delta_{ab}-\Delta_{ab})^2]\big)^{1/2}$ with
    $\Delta_{ab}=\text{RMST}_a-\text{RMST}_b$, the bias of the average effect
    $\mathbb{E}[\hat\Delta_{ab}]-\mathbb{E}[\Delta_{ab}]$, and the share of samples with the
    correct sign. Hidden confounding can affect both bias and PEHE.
    \item \textbf{Policy.} A policy picks the arm with the highest predicted RMST among
    supported arms (abstaining to the assigned arm where none is supported). Value is the
    true RMST of the chosen arm; regret is the oracle value minus it; the hit rate is the share
    of samples where the chosen arm is the true best arm. Reference policies: the assigned
    arm, the best single arm per line (an upper bound for constant policies), and uniform
    random.
    \item \textbf{Factual accuracy (protocol v2).} At the exact requested horizon, write
    $p=S_A(\tau_\ell)$ and $q=\hat S_A(\tau_\ell)$. The primary Brier score averages
    $p(1-q)^2+(1-p)q^2$. We also report Brier scores from uncensored latent outcomes and
    concordance of latent event times with horizon-specific predicted risk. These are
    simulation diagnostics, not estimators available on real censored data. The historical
    IPCW results below used training-cohort censoring weights and must be replaced.
    \item \textbf{References.} The \emph{oracle} (true curves; must score zero error and zero
    regret) and the \emph{naive Kaplan--Meier} (per line and arm, no covariates; the bar a
    covariate model must clear).
\end{itemize}

% -----------------------------------------------------------------------------------------
\subsection{Sweep design}
\label{app:semisynthetic:sweep}

One knob is varied at a time, the others staying at the reference values.
Table~\ref{tab:semisynthetic:sweep} lists the cells; each is replicated three times (data seed
and training seed equal to the replicate index), giving $30$ runs. The hidden-confounder axis
is the one that separates methods that adjust for observed covariates only from methods that
would be unbiased under hidden confounding; the assignment axis is expected to be flat for
outcome-regression estimators, since observed confounding adds variance rather than bias.

\begin{table}[ht]
\centering\small
\caption{Sweep cells. The reference cell is $\gamma=1$, $s=0$, $h=1$; $s=0$ and $h=1$ are
therefore not repeated.}
\label{tab:semisynthetic:sweep}
\begin{tabular}{@{}lll@{}}
\toprule
Knob & Levels & Meaning \\
\midrule
$\gamma$ & $0,\,0.25,\,0.5,\,1,\,1.5$ & assignment strength (observed confounding) \\
$s$ & $0.5,\,1$ (and $0$) & hidden confounder $\to$ assignment \\
$h$ & $0,\,0.5,\,2$ (and $1$) & effect heterogeneity \\
\bottomrule
\end{tabular}
\end{table}

% -----------------------------------------------------------------------------------------
\subsection{Sweep results}
\label{app:semisynthetic:sweep-results}

All $30$ runs of Table~\ref{tab:semisynthetic:sweep} completed and were scored on all four
checkpoint kinds (\texttt{val\_loss}, \texttt{bestCI}, \texttt{bestCALIB}, \texttt{final\_epoch});
Table~\ref{tab:semisynthetic:checkpoint} pools every cell and replicate by kind. \texttt{val\_loss}
--- the real-data recommender's default --- has the highest pooled regret and PEHE of the four;
\texttt{bestCI} is lowest on both and is used for Table~\ref{tab:semisynthetic:sweep-results}.
This was a post-hoc choice using the reported hold-out's oracle metrics, while that hold-out
also supplied early stopping. The following tables are therefore exploratory historical
results, not independent validation. Protocol-v2 aggregation accepts a fixed checkpoint rule
(default \texttt{val\_loss}) and never selects it from the reported oracle regret.

\paragraph{Comparator interpretation.} The omitted Kaplan--Meier effect-error comparison
changes the conclusion: at these historical \texttt{bestCI} checkpoints DynaSurv has higher
mean pair PEHE in nine of ten cells. The reference values are $1.761$ versus $1.705$ months;
at hidden strength 1 they are $2.699$ versus $2.171$; at heterogeneity 0, $2.082$ versus
$1.192$. Only heterogeneity 2 favors DynaSurv ($2.643$ versus $2.773$). Reference policy
regret favors DynaSurv ($0.336$ versus $0.516$ months). The corrected aggregator reports
curve, effect, policy and factual metrics for every comparator on the same support. Ranking
benefit must not be interpreted as accurate individual treatment-effect estimation.

\begin{table}[ht]
\centering\small
\caption{Checkpoint-kind comparison, pooled over all 30 sweep runs. PEHE and $|$bias$|$
(months) are $n$-weighted means over arm pairs and cells; regret (months) and hit rate are
$n$-weighted means of the \texttt{line\_only} policy; factual $C$ is the IPCW concordance index.}
\label{tab:semisynthetic:checkpoint}
\begin{tabular}{@{}lrrrrr@{}}
\toprule
Checkpoint & PEHE & $|$bias$|$ & Regret & Hit rate & Factual $C$ \\
\midrule
\texttt{val\_loss} (real-data default) & 2.37 & 1.51 & 0.42 & 0.62 & 0.666 \\
\texttt{bestCI} & 1.95 & 1.14 & 0.36 & 0.64 & 0.675 \\
\texttt{bestCALIB} & 2.43 & 1.35 & 0.37 & 0.58 & 0.665 \\
\texttt{final\_epoch} & 2.23 & 1.16 & 0.40 & 0.52 & 0.663 \\
\bottomrule
\end{tabular}
\end{table}

\begin{table}[ht]
\centering\small
\caption{Sweep results at the \texttt{bestCI} checkpoint, mean $\pm$ std over 3 replicates.
PEHE and $|$bias$|$ (months) pool over lines and arm pairs; regret (months, oracle $=0$) and
hit rate use the \texttt{line\_only} policy; $h_r$ is the random policy's regret, an upper
bound on the task's headroom at that cell.}
\label{tab:semisynthetic:sweep-results}
\begin{tabular}{@{}llrrrrr@{}}
\toprule
Axis & Level & PEHE & $|$bias$|$ & Regret & Hit rate & $h_r$ \\
\midrule
$\gamma$ & 0        & $1.27\pm0.15$ & $0.34\pm0.15$ & $0.41\pm0.03$ & $0.57\pm0.05$ & 0.77 \\
$\gamma$ & 0.25      & $1.45\pm0.28$ & $0.70\pm0.40$ & $0.45\pm0.02$ & $0.56\pm0.02$ & 0.78 \\
$\gamma$ & 0.5       & $1.61\pm0.19$ & $0.85\pm0.23$ & $0.39\pm0.04$ & $0.59\pm0.04$ & 0.79 \\
$\gamma$ & 1 (ref.)  & $1.76\pm0.42$ & $0.87\pm0.33$ & $0.34\pm0.04$ & $0.64\pm0.04$ & 0.80 \\
$\gamma$ & 1.5       & $2.33\pm1.10$ & $1.49\pm1.03$ & $0.40\pm0.05$ & $0.62\pm0.02$ & 0.81 \\
$s$ & 0.5            & $2.25\pm0.32$ & $1.49\pm0.24$ & $0.44\pm0.04$ & $0.59\pm0.02$ & 0.80 \\
$s$ & 1               & $2.70\pm0.15$ & $2.14\pm0.16$ & $0.48\pm0.01$ & $0.57\pm0.01$ & 0.80 \\
$h$ & 0               & $2.08\pm0.75$ & $1.28\pm0.34$ & $0.03\pm0.03$ & $0.95\pm0.03$ & 0.46 \\
$h$ & 0.5             & $1.46\pm0.29$ & $0.94\pm0.45$ & $0.15\pm0.08$ & $0.67\pm0.12$ & 0.50 \\
$h$ & 2               & $2.64\pm0.53$ & $1.35\pm0.49$ & $0.52\pm0.07$ & $0.62\pm0.01$ & 1.46 \\
\bottomrule
\end{tabular}
\end{table}

\begin{table}[ht]
\centering\small
\caption{Overlap (effective sample size / $n$ per line, mean over 3 replicates), from the
synced manifests. Heterogeneity does not touch assignment, so it repeats the $\gamma=1$ row
exactly at every level.}
\label{tab:semisynthetic:overlap}
\begin{tabular}{@{}llrrrr@{}}
\toprule
Axis & Level & Line 1 & Line 2 & Line 3 & Line 4 \\
\midrule
$\gamma$ & 0        & 0.566 & 0.750 & 0.440 & 0.239 \\
$\gamma$ & 0.25      & 0.539 & 0.709 & 0.422 & 0.217 \\
$\gamma$ & 0.5       & 0.462 & 0.586 & 0.364 & 0.184 \\
$\gamma$ & 1 (ref.)  & 0.251 & 0.282 & 0.198 & 0.086 \\
$\gamma$ & 1.5       & 0.087 & 0.154 & 0.085 & 0.097 \\
$s$ & 0.5            & 0.197 & 0.275 & 0.140 & 0.061 \\
$s$ & 1               & 0.156 & 0.199 & 0.109 & 0.098 \\
$h$ & any level        & 0.251 & 0.282 & 0.198 & 0.086 \\
\bottomrule
\end{tabular}
\end{table}

\paragraph{Reading the three axes.} $\gamma$ (assignment on observed covariates) roughly
doubles PEHE and the average-effect bias from $\gamma=0$ to $\gamma=1.5$, but the model's own
regret is flat within noise (0.34--0.45) while the random-policy headroom $h_r$ grows only
slightly (0.77 to 0.81): the ranking of arms survives better than the size of the effect does.
This is consistent with overlap collapse: mean effective sample size / $n$ on line 4 falls from
$0.24$ at $\gamma=0$ to $0.10$ at $\gamma=1.5$ (Table~\ref{tab:semisynthetic:overlap}), so
precision, not direction, is what degrades. Three replicates cannot separate true bias from
finite-sample noise on this axis; the sweep gives a variance estimate, not a significance test.
$s$ (hidden confounder into assignment) is the sharper effect: bias more than doubles from the
$s=0$ reference ($0.87$) to $s=1$ ($2.14$) and regret rises by $40\%$ even though its overlap
loss is comparable to $\gamma$'s (line-4 ESS/$n$ falls to $0.06$--$0.10$) -- the pattern expected
of a confounder the model cannot see, not merely thinner support. $h$ (heterogeneity) mainly
changes the size of the task rather than the model's absolute accuracy: it does not change
assignment at all (identical overlap to the $\gamma=1$ reference row, Table~\ref{tab:semisynthetic:overlap}),
so any shift in the metrics below comes from the outcome model alone. At $h=0$ every line has
a single best arm, so both the model and a constant policy score near zero regret ($h_r=0.46$
is administrative-censoring headroom, not heterogeneity); as $h$ grows to $2$, $h_r$ triples to
$1.46$ and the model's regret grows with it (to $0.52$) while staying well under the random
policy's ceiling and the hit rate falls from $0.95$ to $0.62$ -- picking the right arm gets
harder as there are more genuinely different arms to pick between.

\paragraph{What this does not show.} The 30-run sweep uses one training configuration
(Section~\ref{app:semisynthetic:training}) and the model config predating this benchmark; it is
not re-tuned per cell. It has no Cox or DeepSurv baseline, so it cannot yet reproduce the
observed-vs-hidden-confounding bias comparison of Section~\ref{sec:counterfactual} with this
generator. Achieved censoring event rates matched their per-line targets to within $0.0007$ in
every one of the 30 synced manifests, confirming the calibration of
Section~\ref{app:semisynthetic:censoring} generalizes across the sweep.

% -----------------------------------------------------------------------------------------
\subsection{Verification of the generator}
\label{app:semisynthetic:verification}

The following were checked on the reference replicate.
\begin{itemize}
    \item Mean propensity per line equals the target to $10^{-6}$ for
    $\gamma\in\{0,0.5,1,2,4\}$ and $s\in\{0,2\}$; sampled arm counts pass a chi-square test;
    at $\gamma=s=0$ the propensity is constant within a line.
    \item A multinomial regression of the sampled arm on the drivers, pooled over lines with
    line indicators, recovers the assignment coefficients (correlation $0.976$, all signs
    correct). On line 1 alone it does not ($0.58$), because the previous-line drivers are
    constant there.
    \item Simulated Kaplan--Meier curves match the real ones (within about $0.04$ at 6, 12,
    24 and 36 months, medians within $2\%$); observed curves after censoring within about
    $0.02$; simulated event rates equal the real rates.
    \item Closed-form RMST agrees with Monte Carlo (bias $-0.002$ months, standardised
    differences of standard deviation $1.01$); at $h=0$, $\tau_a=0$ all arms have equal $\eta$.
    \item Expansion: the unmasked rows per line equal the real patients per line; history rows
    equal the real records; no hidden variable and no propensity in the model-facing files.
    \item Reproducibility: an identical rerun is bit-identical; changing $\gamma$ keeps $u$
    and changes $13\%$ of the arms; a new replicate changes $u$.
    \item No patient appears on both sides of a temporal, random or cross-validation split.
\end{itemize}

% -----------------------------------------------------------------------------------------
\subsection{Limitations}
\label{app:semisynthetic:limitations}

\begin{enumerate}
    \item Coefficients are hand-set; results describe the difficulty of this generator, not
    real effect sizes.
    \item One Weibull shape per line, a shared coupling across arms, non-informative censoring,
    no outcome feedback across lines, and no treatment change within a line.
    \item Only $16$ covariates affect the outcome; the model sees many more real columns, which
    are non-causal here by construction.
    \item Arms are thin on late lines (\arm{ET alone} and \arm{ET+anti-CDK} on line 4): their
    counterfactuals are extrapolation, and unsupported arms are excluded from the metrics.
    \item Calendar time is an instrument and the hold-out enters late: an intentional shift
    that the reported errors include.
    \item The covariate standardiser is fitted on all samples, hold-out included (a property of
    the existing data module).
    \item Training settings were tuned on real data and chosen by hand here; the benchmark
    scores the model \emph{as trained}, not the best model it could be.
\end{enumerate}
